\documentclass{article}

%% AMS packages
\usepackage{amsfonts,amsmath,amstext,amsbsy,amssymb}
\usepackage{amsopn,amsthm,amscd,amsxtra}

\usepackage{dsfont}

%% Graphics
\usepackage{graphicx}

%% Figures
% \usepackage{epsfig} \input{xy} \xyoption{all}

 \usepackage{url}


%% LATeX packages
\usepackage{latexsym,mathrsfs,stmaryrd}


\usepackage{thmtools}
\usepackage{setspace}
\usepackage{geometry}
\usepackage{float}
\usepackage{hyperref}
\usepackage[utf8]{inputenc}
\usepackage[english]{babel}
\usepackage{framed}
\usepackage[dvipsnames]{xcolor}
\usepackage{tcolorbox}


%% Special arrows
\usepackage{mathabx} \newcommand{\carr}{\righttoleftarrow}
\newcommand{\wto}{\rightharpoonup}

%% Hyperlinks setup
\usepackage{hyperref}
\RequirePackage[dvipsnames]{xcolor} % [dvipsnames]
\definecolor{halfgray}{gray}{0.55}
% chapter numbers will be semi transparent .5 .55 .6 .0
\definecolor{webgreen}{rgb}{0,0.5,0}
\definecolor{webbrown}{rgb}{.6,0,0} \hypersetup{%
  colorlinks=true, linktocpage=true, pdfstartpage=3,
  pdfstartview=FitV,%
  breaklinks=true, pdfpagemode=UseNone, pageanchor=true,
  pdfpagemode=UseOutlines,%
  plainpages=false, bookmarksnumbered, bookmarksopen=true,
  bookmarksopenlevel=1,%
  hypertexnames=true,
  pdfhighlight=/O,%hyperfootnotes=true,%nesting=true,%frenchlinks,%
  urlcolor=webbrown, linkcolor=RoyalBlue,
  citecolor=webgreen, %pagecolor=RoyalBlue,%
  % uncomment the following line if you want to have black links
  % (e.g., for printing)
  % urlcolor=Black, linkcolor=Black,
  % citecolor=Black, %pagecolor=Black,%
  pdftitle={},%
  pdfauthor={},%
  pdfsubject={2000 Mathematical Subject Classification: Primary:},%
  pdfkeywords={},%
  pdfcreator={pdfLaTeX},%
  pdfproducer={LaTeX with hyperref}%
}


\newcommand{\db}[1]{\llbracket #1\rrbracket}



%%%%%%%%%%%%%%%%%%%%%%%%%%%%%%

%% Portuguese
% \usepackage[brazilian]{babel}
\usepackage{ae}

%%%%%%%%%%%%%%%%%%%%%%%%%%%%%%

% Bracket-like operations
\newcommand{\floor}[1]{\left\lfloor{#1}\right\rfloor}
\newcommand{\ceil}[1]{\left\lceil{#1}\right\rceil}
\newcommand{\abs}[1]{\left\lvert{#1}\right\rvert}
\newcommand{\norm}[1]{\left\|{#1}\right\|}
\newcommand{\gen}[1]{\left\langle{#1}\right\rangle}
\newcommand{\scprod}[2]{\left\langle{#1},{#2}\right\rangle}

% Typefaces
\newcommand{\bb}{\mathbb} \newcommand{\mbf}{\mathbf}
\newcommand{\mc}{\mathcal} \newcommand{\ms}{\mathscr}
\newcommand{\ds}{\mathds}

% Usual spaces
\newcommand{\R}{\mathbb{R}}\newcommand{\N}{\mathbb{N}}
\newcommand{\Z}{\mathbb{Z}}\newcommand{\Q}{\mathbb{Q}}
\newcommand{\T}{\mathbb{T}}\newcommand{\C}{\mathbb{C}}
\newcommand{\D}{\mathbb{D}}\newcommand{\A}{\mathbb{A}}

%% Abbreviations
\newcommand{\cf}{c.f.\ } \newcommand{\ie}{i.e.{}}
\newcommand{\eg}{e.g.\ } \newcommand{\etal}{\textit{et al\ }}


%% Diff-like operators
\newcommand{\Hom}{\mathrm{Hom}}
\newcommand{\Diff}[2]{\mathrm{Diff}_{#1}^{#2}}
\newcommand{\Homeo}[1]{\mathrm{Homeo}_{#1}}
\newcommand{\Dis}[1]{\mathcal{D}^\prime_{#1}}


%% Measures and differentials
\newcommand{\Leb}{\mathrm{Leb}} \newcommand{\dd}{\:\mathrm{d}}

%% Some operators
\DeclareMathOperator{\M}{\mathfrak{M}} \DeclareMathOperator{\Lip}{Lip}
\DeclareMathOperator{\Per}{Per} \DeclareMathOperator{\Fix}{Fix}
\DeclareMathOperator{\SL}{SL} \DeclareMathOperator{\GL}{GL}
\newcommand{\gl}{\mathfrak{gl}}
\DeclareMathOperator{\spec}{spec}
\DeclareMathOperator{\rank}{rank} \DeclareMathOperator{\diver}{div}
\DeclareMathOperator{\cl}{cl} \DeclareMathOperator{\sign}{sign}
\DeclareMathOperator{\diam}{diam}

\newcommand{\cc}{\mathrm{cc}} % connected components

\newcommand{\pr}[1]{\mathrm{pr}_{#1}} % projections


%% For marginal notes
\newcommand{\marginal}[1]{\marginpar{{\scriptsize {#1}}}}

%% Topological dynamics on tori
\newcommand{\Symp}[1]{\mathrm{Symp}_{#1} (\mathbb{T}^2)}
\newcommand{\Ham}{\mathrm{Ham} (\mathbb{T}^2)}
\newcommand{\Thomeo}{\widetilde{\mathrm{Homeo}_0}}
\newcommand{\Tsymp}{\widetilde{\mathrm{Symp}_0} (\mathbb{T}^2)}
\newcommand{\Tham}{\widetilde{\mathrm{Ham}} (\mathbb{T}^2)}
\newcommand{\wh}{\widehat}


%% Especial notation for this article

\newcommand{\B}{\mathscr{B}}
\newcommand{\Bs}{\mathcal{S}}
\newcommand{\eps}{\varepsilon}
\DeclareMathOperator{\MS}{MS}
\newcommand{\Or}{\mathcal{O}}
\newcommand{\Rec}{\mathrm{Rec}}
\newcommand{\CR}{\mathscr{R{}}}
\newcommand{\RP}{\mathbb{RP}}
\DeclareMathOperator{\osc}{osc}


%%% Settings of original template

\colorlet{LightGray}{White!90!Periwinkle}
\colorlet{LightOrange}{Orange!15}
\colorlet{LightGreen}{Green!15}
\colorlet{LightYellow}{Yellow!15}
\colorlet{LightBlue}{Blue!15}

\newcommand{\HRule}[1]{\rule{\linewidth}{#1}}

\declaretheoremstyle[name=Theorem,]{thmsty}
\declaretheorem[style=thmsty,numberwithin=section]{theorem}
\tcolorboxenvironment{theorem}{colback=LightOrange}

\declaretheoremstyle[name=Corollary,]{corsty}
\declaretheorem[style=corsty,numberlike=theorem]{corollary}
\tcolorboxenvironment{corollary}{colback=LightOrange}


\declaretheoremstyle[name=Proposition,]{prosty}
\declaretheorem[style=prosty,numberlike=theorem]{proposition}
\tcolorboxenvironment{proposition}{colback=LightYellow}

\declaretheoremstyle[name=Lemma,]{lemsty}
\declaretheorem[style=lemsty,numberlike=theorem]{lemma}
\tcolorboxenvironment{lemma}{colback=LightGreen}

\declaretheoremstyle[name=Definition,]{defsty}
\declaretheorem[style=defsty,numberwithin=section]{definition}
\tcolorboxenvironment{definition}{colback=LightBlue}


\declaretheoremstyle[name=Problem,]{probsty}
\declaretheorem[style=probsty,numberwithin=section]{problem}
\tcolorboxenvironment{problem}{colback=LightGray}

\declaretheoremstyle[name=Example,]{examsty}
\declaretheorem[style=exasty,numberwithin=section]{example}
\tcolorboxenvironment{example}{colback=LightBlue}


\declaretheoremstyle[name=Exercise,]{exesty}
\declaretheorem[style=exesty,numberlike=problem]{exercise}
\tcolorboxenvironment{exercise}{colback=LightGray}



\setstretch{1.2}
\geometry{
    textheight=9in,
    textwidth=5.5in,
    top=1in,
    headheight=12pt,
    headsep=25pt,
    footskip=30pt
}


 %% Fancy indexing in enumerations
\usepackage{enumerate}

%% Figures
% \usepackage{epsfig} \input{xy} \xyoption{all}

\usepackage{url}



\setstretch{1.2}
\geometry{
    textheight=9in,
    textwidth=5.5in,
    top=1in,
    headheight=12pt,
    headsep=25pt,
    footskip=30pt
}

% ------------------------------------------------------------------------------

\begin{document}

% ------------------------------------------------------------------------------
% Cover Page and ToC
% ------------------------------------------------------------------------------

\title{ \normalsize \textsc{Doutorado em Matemática - IMPA}
        \\ [2.0cm]
        \HRule{1.5pt} \\
        \LARGE \textbf{\uppercase{Lecture notes on Dynamics I}
        \HRule{2.0pt} \\ [0.6cm] \LARGE{} \vspace*{10\baselineskip}}}
\date{}
\author{Alejandro Kocsard \\
        IMPA\\
        \today}

\maketitle
\newpage

\tableofcontents
\newpage

% ------------------------------------------------------------------------------


\section{Introduction}\label{sec:introduction}

In this course we shall deal with \emph{autonomous deterministic dynamical systems.} They can be used to represent or modelize any phenomenon that changes along the time in such a way that the current state of the system completely determines the future evolution and the law that drives this evolution does not change along the time. Mathematically, an \textbf{autonomous deterministic dynamical system} (a \textbf{dynamical system}, for short) is given by a set $M$, usually called the \textbf{phase space,} which represents the set of possible states of the system, a subsemigroup $\T\subset\R$ that represents the \textbf{time set,} and \emph{``a law''} $\Phi\colon\T\times M\to M$ that determines the evolution of the system in such a way that given the current state of our system $x_{0}\in M$, at time $t\in\T$ the corresponding state will be $x_{t}=\Phi(x_{0},t)$. The autonomy of the system is reflected in the so called \textbf{flow} (or \textbf{semi-flow}) condition function $\Phi$ satisfies:
\begin{equation}\label{eq:flow-condition}
  \Phi(t+s,x)=\Phi\big(t,\Phi(s,x)\big) = \Phi\big(t,\Phi(s,x)\big), \quad\forall t,s\in\T,\ \forall x\in M.
\end{equation}

We will consider two different kinds of dynamical systems: \textbf{discrete time} and \textbf{continuous time} ones, \ie~where $\T=\N_{0}$ and $\T=\R_{\geq 0}$, respectively. Moreover, we will also consider \textbf{invertible systems} where the current state does not just determine the future but the past as well. When dealing with invertible discrete dynamical systems, the time set will be $\T=\Z$, and in the case of invertible continuous time ones, we will have $\T=\R$.

The main goal of Dynamical Systems consists in describing the \textbf{orbits of the systems,} \ie~the sets
\begin{displaymath}
  \Or_{\Phi}(x):=\left\{\Phi(t,x) : t\in\T\right\}, \quad\forall x\in M.
\end{displaymath}


\begin{exercise}[Discrete time generator]
  Given a discrete time dynamical system $\Phi\colon\T\times M\to M$, show that $\Phi$ satisfies the flow condition~\eqref{eq:flow-condition} if and only if there is $f\colon M\carr$ such that
  \begin{displaymath}
    \Phi(n,x)=f^{n}(x), \quad\forall x\in M,\ \forall n\in\N_{0},
  \end{displaymath}
  where the iteration of $f$ is inductively defined by $f^{0}:=id_{M}$ and $f^{n+1}=f\circ f^{n}$, for every $n\geq 0$. Show that $\Phi$ is invertible if and only if $f$ is bijective and, in such a case, one can define $f^{-n}:={\big(f^{-1}\big)}^{n}$, for all $n\geq 0$.
\end{exercise}

On the other hand, one can find \emph{infinitesimal generators} for continuous time smooth dynamical systems:

\begin{exercise}[Continuous time generator]
  Let $M$ be a differentiable closed manifold (\ie~compact and boundaryless) and $\Phi\colon\R\times M\to M$ a continuous time dynamical system (that satisfies the flow condition~\ref{eq:flow-condition}). If  we assume $\Phi$ is a $C^{2}$ map and we define $X(p):=\partial_{t}\Phi(t,p)\big|_{t=0}\in T_{p}M$, for every $p\in M$, then the function $\Phi(\cdot,x_{0})\colon\R\to M$ is the unique solution of the \emph{initial value problem}
  \begin{displaymath}
    \dot x(t)=X\big(x(t)\big), \quad x(0)=x_{0},
  \end{displaymath}
  for every $x_{0}\in M$. In other words, the vector field $X$ is the infinitesimal generator of the flow $\Phi$.
\end{exercise}

As previously mentioned, the primary objective of Dynamical Systems is to understand the orbits of various systems, particularly their asymptotic behavior as time goes to infinity. From a historical perspective, there are three pivotal moments in the development of this mathematical theory. The first can be attributed to Isaac Newton, who laid the groundwork for Calculus and employed its principles to solve differential equations. During this era, the approach to studying a specific continuous-time dynamical system involved solving the corresponding differential equation, followed by an analysis of the obtained solutions.

The so-called \emph{``restricted two-body problem''} is a remarkable question that was answered at that time:

\begin{problem}
  The \emph{``restricted two-body problem''} refers to a simplified model used to analyze the motion of two celestial bodies under the influence of their mutual gravitational attraction, where one body is significantly more massive than the other. In this framework, the more massive body is assumed to be stationary, while the second body is treated as a point mass that moves in the gravitational field created by the first one. This approach allows for the derivation of equations of motion for the second body without needing to account for its gravitational effects on the primary, simplifying the analysis. This model is particularly useful in celestial mechanics, where it can be applied to scenarios such as a satellite orbiting a planet or a comet passing by a star.

  More precisely, the \emph{restricted two-body problem} consider in describing the orbit of massive point $p$, whose location in the space at time $t\in\R$ is given by $x(t)\in\R^{3}$, moving in a central field force in $\R^{3}$ that is generated by a massive point of mass $M>0$ which is fixed at the origin. According to Newton's laws, the trajectory $x(t)\in\R^{3}$ is the solution of the following initial value problem:
  \begin{equation}
    \label{eq:restricted-2-body-problem}
    \left\{
      \begin{aligned}
        \ddot{x}(t) &= -\frac{GM}{\norm{x}^{3}}x(t),\\
        x(0)&=x_{0},\\
        \dot{x}(0)&=v_{0},
      \end{aligned}
    \right.
  \end{equation}
  where $G>0$ is the universal constant of gravitation,  $x_{0}\in\R^{3}$ and $v_{0}\in\R^{3}$ are the initial position and velocity, respectively.

  Using classical \emph{quadrature techniques} one can show that the solutions $x(t)$ of~\eqref{eq:restricted-2-body-problem} are either straight lines through the origin (when $x_{0}$ and $v_{0}$ are collinear vectors in $\R^{3}$), or quadratic curves (\ie~ellipses, parabolas or hyperbolas).
\end{problem}

The second pivotal moment in the development of Dynamical Systems can be attributed to Henri Poincaré. At the end of XIXth century it was already pretty clear that there were many ordinary differential equations that could not be solved with quadrature methods. A very simple and well-known example is the following one:

\begin{example}
  There is no \emph{elementary (smooth) function} $f\colon\R\to\R$ such that
  \begin{displaymath}
    f'(x)=e^{-x^{2}}, \quad\forall x\in\R.
  \end{displaymath}
\end{example}

Another example, which was fundamental in development of so called \emph{qualitative theory of differential equations} is the \emph{$3$-body problem:}

\begin{problem}[$n$-body problem]\label{prob:n-body-problem}
  The \emph{$n$-body problem} refers to the challenge of predicting the individual motions of a group of $n$ celestial objects that interact with each other through gravitational forces. The complexity arises from the fact that each body exerts a gravitational force on every other body, leading to a system of coupled differential equations. If $x_{i}(t)\in\R^{3}$ represents the position of the $i$-th object at time $t\in\R$ and $m_{i}>0$ denotes its mass, then Newton's second law determines that these quantities are related by the following system of ordinary differential equations
  \begin{equation}
    \label{eq:n-body-problem}
    \left\{
      \begin{aligned}
        \ddot x_{i}(t) &= \sum_{j\neq i} G\frac{m_{j}}{\norm{x_{i}(t)-x_{j}(t)}^{3}}\big(x_{i}(t)-x_{j}(t)\big), \\
        x_{i}(0)&=x_{i,0}, \\
        \dot x_{i}(0)&=v_{i,0},
      \end{aligned}
    \right.
  \end{equation}
  for $1\leq i\leq n$, and where $x_{i,0},v_{i,0}\in\R^{3}$ denotes the position and the velocity of the $i$-th object at time $t=0$, respectively.

  It is pretty remarkable that, in general, the system of ODEs given by~\eqref{eq:n-body-problem} cannot be solved by ``quadrature methods'' when $n\geq 3$. In fact, their solutions are not elementary functions in general.
\end{problem}

The $n$-body problem, as outlined in Problem~\ref{prob:n-body-problem}, served as a primary motivation for Henri Poincaré in developing the Geometric Theory of Differential Equations. Poincaré posited that if explicit solutions to a given differential equation cannot be readily obtained, one should at least strive to describe some interesting geometric properties of those solutions.

Indeed, many questions regarding the solutions of differential equations are more qualitative than quantitative, or more geometric than analytic. For instance, a fundamental qualitative inquiry concerning the system of ordinary differential equations in~\eqref{eq:n-body-problem} is whether every solution is periodic. In some cases, even when solutions can be expressed using elementary functions, addressing certain qualitative questions based solely on the explicit formulas can prove to be quite challenging.

The third pivotal moment in the evolution of Dynamical Systems can be attributed to the schools of Steve Smale in the US and Aleksandr Andronov in the USSR. At this juncture, it became evident that the geometric properties of orbits, trajectories, or solutions to ODEs could change dramatically with variations in system parameters. For example, in Problem~\ref{prob:n-body-problem}, there are $3n$ parameters—comprising the $n$ masses, $n$ initial positions, and $n$ initial velocities. In certain instances, even minor changes to these parameters can lead to completely different geometric properties for the respective solutions. Thus, it becomes crucial to identify geometric properties of orbits that are robust or persistent under small perturbations of system parameters.

These three historical stages in the development of dynamical systems are cleverly (and provocatively) summarized by Étienne Ghys in his laudation for Artur Avila’s Fields Medal: In Newton's time, you were given an ODE and tasked with finding its solutions. In Poincaré's era, you were given an ODE and asked to elucidate interesting geometric properties of its solutions. In Smale's time, you were given no ODE at all and challenged to say something interesting!


\section{Some examples of Dynamical Systems results}\label{sec:some-exampl-dynamical-sys-res}

A very first example of a result in Dynamical Systems might be the following classical

\begin{theorem}[Banach fixed point theorem]\label{thm:Banach-fixed-point-thm}
  Let $(M,d)$ be a complete metric space and $f\colon M\carr$ be \textbf{uniform contraction,} \ie~$f$ is a Lipschitz function and its \textbf{Lipschitz constant}
  \begin{displaymath}
    \lambda=\Lip(f):=\sup_{x\neq y}\frac{d\big(f(x),f(y)\big)}{d(x,y)}< 1.
  \end{displaymath}

  Then there exists a unique point $p\in M$ such that $f(p)=p$ and it holds
  \begin{displaymath}
    \lim_{n\to+\infty} f^{n}(x)=p, \quad\forall x\in M.
  \end{displaymath}
\end{theorem}

\begin{proof}
  Given any $x\in M$, one can show by induction on $n$ that
  \begin{displaymath}
    d\big(f^{n}(x),f^{n+1}(x)\big)\leq \lambda^{n}d(x,f(x)), \quad\forall n\in\N_{0}.
  \end{displaymath}
  So, for every $m,n\in\N$ it holds
  \begin{displaymath}
    \begin{split}
      d\big(f^{m}(x),f^{m+n}(x)\big)&\leq \sum_{i=0}^{n-1} d\big(f^{m+i}(x),f^{m+i+1}(x)\big) \\
      &\leq \sum_{i=0}^{n-1}\lambda^{m+i}d\big(x,f(x)\big) \leq \frac{\lambda^{m}}{1-\lambda} d\big(x,f(x)\big).
    \end{split}
  \end{displaymath}
  Hence, the sequence ${\big(f^{n}(x)\big)}_{n\geq 0}$ is Cauchy, and since $M$ is complete, it converges to a point $p\in M$. By continuity of $f$, $p$ is a fixed point of $f$. On the other hand, since $f$ is a contraction, it has at most one fixed point.
\end{proof}

\begin{exercise}[Perturbations of contractions]
  Let $f\colon\R^{d}\carr$ be a uniform contraction and $r>0$ be arbitrary. We define the following balls
  \begin{displaymath}
    B_{r}^{0}(f):=\left\{g\in C^{0}(\R^{d},\R^{d}) : \norm{f-g}_{C^{0}}:=\sup_{x\in\R^{d}}\norm{f(x)-g(x)}<r\right\},
  \end{displaymath}
  and
  \begin{displaymath}
    B_{r}^{L}(f):=\left\{g\in C^{0}(\R^{d},\R^{d}) : \norm{f-g}_{L}:=\norm{f-g}_{C^{0}} + \Lip(f-g)<r\right\}.
  \end{displaymath}

  Show that given any $r>0$, there is a $g\in B_{r}^{0}(f)$ which is not a contraction and such that $\Fix(g):=\{x\in\R^{d} : g(x)=x\}$ is an infinite set.

  On the other hand, show that given the uniform contraction $f$, there exists $r_{0}>0$ such that every $g\in B_{r_{0}}^{L}(f)$ is a uniform contraction, and hence it exhibits a unique fixed point $p_{g}\in\R^{d}$. Finally, prove that given any $\eps>0$, there exists $\delta\in(0,r_{0})$ such that
  \begin{displaymath}
    \norm{p_{f}-p_{g}}<\eps, \quad\forall g\in B_{\delta}^{L}(f).
  \end{displaymath}
\end{exercise}

Another important example is the following one:

\begin{example}[Gradient flows]\label{exam:gradient-flows}
  Let $M$ be the unit sphere $\bb{S}^{2}:=\{v\in\R^{3} : \norm{v}=1\}$ and consider the function $L\colon\bb{S}^{2}\to\R$ given by
  \begin{displaymath}
    L(x,y,z):=-z, \quad\forall (x,y,z)\in\bb{S}^{2},
  \end{displaymath}
  and the flow $\Phi\colon\R\times\bb{S}^{2}\to\bb{S}^{2}$ generated by the vector field $X(v):=-\nabla L(v)$, for every $v\in\bb{S}^{2}$ and where $\nabla$ denotes the gradient operator on $\bb{S}^{2}$.

  One can easily show that $(0,0,-1)$ and $(0,0,1)$ are the only singularities of $X$ (\ie~the critical points of $L$), or equivalently, the fixed points of $\Phi$, and it holds
  \begin{displaymath}
    \Phi^{t}(v)\to (0,0,-1), \quad\text{and}\quad \Phi^{-t}(v)\to (0,0,1),
  \end{displaymath}
  as $t\to+\infty$, for every $v\in\bb{S}^{2}\setminus\{(0,0,\pm 1)\}$.

  Generalizing this example, one can consider an arbitrary compact differentiable manifold $M$ endowed with a Riemannian structure $\langle\cdot,\cdot\rangle$ and a smooth function $L\colon M\to\R$. Then we can define the gradient of $L$ as the only vector field $\nabla L$ such that
  \begin{displaymath}
    YL = dL(Y) =\langle\nabla L, Y\rangle,
  \end{displaymath}
  for every vector field $Y$ on $M$.

  So, if we consider the flow $\Phi\colon\R\times M\to M$ induced by the vector field $X:=-\nabla L$, one can easily check that
  \begin{displaymath}
    \partial_{t} (L\circ\Phi^{t})(p)\big|_{t=0}=-\norm{\nabla L(p)}^{2} = -\langle\nabla L(p),\nabla L(p)\rangle,\quad\forall p\in M.
  \end{displaymath}

  That means, given any point $p\in M$ the map $\R\ni t\mapsto L\circ\Phi^{t}(p)\in\R$ is not increasing, and it is strictly decreasing if and only if $p$ is a regular point for $L$ is a regular point of $L$, \ie~$dL_{p}\colon T_{p}M\to\R$ is not identically zero.

  Moreover, if $L$ has a finite number of critical points (\ie~points where the derivative $dL$ vanishes), then one can show that for every $p\in M$, there exist critical points of $L$ $\alpha_{p},\omega_{p}\in M$ such that
  \begin{displaymath}
    \Phi^{-t}(p)\to\alpha_{p}, \quad\text\quad \Phi^{t}(p)\to\omega_{p},
  \end{displaymath}
  as $t\to+\infty$.
\end{example}


\section{The Fundamental Theorem of Dynamical Systems}\label{sec:fundamental-thm-dyn-sys}

Motivated by Example~\ref{exam:gradient-flows}, we would like to find functions decreasing along the orbits of arbitrary topological dynamical systems.

\begin{definition}[Lyapunov function]
  Given a continuous dynamical systems $f\colon M\carr$ on a metric space $M$, a \emph{Lyapunov function} for $f$ is a continuous function $L\colon M\to\R$ such that
  \begin{displaymath}
    L\big(f(x)\big)\leq L(x), \quad\forall x\in M.
  \end{displaymath}
\end{definition}

Of course, according to this definition every constant function is a Lyapunov one for every topological dynamical system, which does not seem to be very interesting. So, we are going to look for the ``best possible'' Lyapunov function.

\subsection{Different forms of recurrence}\label{sec:forms-of-recurrence}

It would be desirable to find a Lyapunov function such that $L\big(f(x)\big)<L(x)$ for every $x\in M$. However this is clearly impossible when $f$ has, for instance, a \textbf{fixed point,} \ie~a point $p\in M$ such that $f(p)=p$. The set of fixed points of $f$ will be denoted by $\Fix(f)$.

More generally, we write $p\in\Per(f)$ and say that $p$ is a \textbf{periodic point} when there is $n\geq 1$ such that $f^{n}(p)=p$. We define $\Per(f):=\bigcup_{n\geq 1}\Fix(f^{n})$.

Another form of recurrence that represents an obstruction for the existence of a Lyapunov functions satisfying the strict inequality for every $p\in M$ is the following: we say that a point $p\in M$ is \textbf{positively recurrent} for $f$ when there exists strictly increasing sequence of natural numbers $n_{j}\uparrow+\infty$ such that $f^{n_{j}}(p)\to p$, as $n_{j}\to+\infty$. The set of positively recurrent points of $f$ shall be denoted by $\Rec^+(f)$. When $f\colon M\carr$ is bijective, one says that $p\in M$ is \textbf{negatively recurrent} for $f$ when it is positively recurrent for $f^{-1}$. The set of negatively recurrent points will be denoted by $\Rec^-(f)$.

Given a point $p\in M$, we define its \textbf{$\omega$-limit set} by
\begin{displaymath}
  \omega_{f}(p):=\left\{x\in M : \exists n_{j}\uparrow\infty,\ f^{n_{j}}(p)\to x,\ \text{as } n_{j}\to+\infty\right\},
\end{displaymath}
and the \textbf{positive limit set} of $f$ by
\begin{displaymath}
  L^+(f):=\bigcup_{p\in M} \omega_{f}(p).
\end{displaymath}

When $f$ is invertible, we also define the \textbf{$\alpha$-limit set} of $p\in M$ by $\alpha_{f}(p):=\omega_{f^{-1}}(p)$, and the so called \textbf{limit set} of $f$ by
\begin{displaymath}
  L(f):=\bigcup_{p\in M} \alpha_{f}(p)\cup \omega_{f}(p).
\end{displaymath}

We say that a point $p\in M$ is \textbf{non-wandering} when for every neighborhood $U$ of $p$, there exists $n\geq 1$ such that $U\cap f^{-n}(U)\neq\emptyset$. The set of non-wandering points for $f$ is usually denoted by $\Omega(f)$.

We say that a subset $K\subset M$ is \textbf{$f$-invariant} when $f(K)=K$, and it is said to be \textbf{totally $f$-invariant} when $f(K)=K=f^{-1}(K)$. Observe that if $f\colon M\carr$ is bijective, then $K\subset M$ is totally $f$-invariant if and only if it is $f$-invariant.

\begin{exercise}
  \label{ex:inclusions-fix-per-rec-limit-sets-non-wandering}
  Let $M$ denote a compact metric space and $f\colon M\carr$ be a homeomorphism. Show that
  \begin{displaymath}
    \Fix(f)\subset\Per(f)\subset\Rec^+(f)\subset L(f)\subset\Omega(f),
  \end{displaymath}
  and all these sets are $f$-invariant.

  \begin{enumerate}
    \item Show that $\Fix(f)$ and $\Omega(f)$ are compact.
    \item Show that $\Rec^+(f)\cap\Rec^-(f)\neq\emptyset$. (Hint: show that there exists a non-empty compact subset $K\subset M$ which is \textbf{minimal} for $f$, \ie~$K$ is $f$-invariant and it does not contain any proper compact $f$-invariant subset).
    \item Show that if $L\colon M\to\R$ is a Lyapunov function for $f$ and $p\in\Omega(f)$, then it holds $L\big(f(p)\big)=L(p)$.
  \end{enumerate}
\end{exercise}

Finally we introduce the notion of \emph{chain recurrence.} Given a homeomorphism $f\colon M\carr$ on the compact metric space $(M,d)$ and an arbitrary $\eps>0$, a sequence ${(x_{n})}_{n\in\Z}$ of points of $M$ is said to be an \textbf{$\eps$-chain} or \textbf{$\eps$-pseudo-orbit} when
\begin{displaymath}
  d\big(f(x_{i}),x_{i+1}\big)<\eps,\quad\forall n\in\Z.
\end{displaymath}

A point $p\in M$ is said to be \textbf{chain recurrent} when for every $\eps>0$ there exists a periodic $\eps$-chain ${(x_{n})}_{n\in\Z}$ with $p=x_{0}$. The set of chain recurrent points is denoted by $\CR(f)$.


\begin{exercise}
  Show that $\CR(f)$ is a compact $f$-invariant set and $\Omega(f)=\Omega(f^{-1})\subset\CR(f)=\CR(f^{-1})$.
\end{exercise}

Given any $p\in M$ and any $\eps>0$, we define the \textbf{$\eps$-chainable set from $p$} as the set
\begin{displaymath}
  \CR_{f}(p,\eps):=\left\{x\in M : \exists {(x_{n})}_{n\in\Z}\ \eps\text{-chain, } \exists k\in\N, x_{0}=p,\ x_{k}=x\right\},
\end{displaymath}
and the \textbf{chainable set from $p$} by
\begin{displaymath}
  \CR_{f}(p):=\bigcap_{\eps>0}\CR_{f}(p,\eps).
\end{displaymath}
So, we have $p\in\CR(f)$ if and only if $p\in\CR_{f}(p)$.

There is a fundamental and elementary notion that will play a crucial role in the understanding of the orbit structure of a given dynamical system, which is the concept of \emph{trapping region.} We say that a subset $U\subset M$ is a \textbf{trapping region} for $f$ when it is open and $f\big(\overline{U}\big)\subset U$.

Observe that if $U\subset M$ is a trapping region for $f$, then the set $M\setminus\overline{U}$ is a trapping region for $f^{-1}$; and reciprocally, if $M\setminus\overline{U}$ is a trapping for $f^{-1}$, then $\mathrm{int}(\overline{U})$ is a trapping region for $f$.

Given a trapping region $U$ for $f$, we can define the \textbf{dual pair of attractor/repellor} associated to $U$ as the pair of compact $f$-invariant sets
\begin{displaymath}
  A_{U}^+:=\bigcap_{n\geq 0} f^{n}\big(\overline{U}\big), \quad\text{and}\quad A_{U}^-:=\bigcap_{n\geq 0} f^{-n}\big(M\setminus U\big).
\end{displaymath}

\begin{exercise}
  Show that $A_U^+$ and $A_U^-$ are compact, $f$-invariant and contained in $U$ and $M\setminus\overline{U}$, respectively. Moreover, they are non-empty if and only if $U$ and $M\setminus\overline{U}$ are non-empty.
\end{exercise}

We have the following fundamental result relating chain recurrence and trapping regions:
\begin{proposition}
  \label{prop:chain-recurrence-trapping-regions}
  Let $f\colon M\carr$ be a homeomorphism on the compact metric space $(M,d)$. Then, given any point $p\in M$ and any $\eps>0$, the $\eps$-chainable set from $p$ is a trapping region for $f$.
\end{proposition}

\begin{proof}
  Let us start showing that $\CR_f(p,\eps)$ is open. To do that, let us inductively defined
  \begin{displaymath}
    U_1:=B_\eps(p)=\{y\in M : d(x,y)<\eps\},
  \end{displaymath}
  and
  \begin{displaymath}
    U_n:=\bigcup_{x\in U_{n-1}} B_\eps\big(f(x)\big), \quad\forall n\geq 2.
  \end{displaymath}

  One can easily check that $U_n$ is open for every $n\geq 1$ and $\CR_f(p,\eps)=\bigcup_{n\geq 1} U_n$. So, $\CR_f(p,\eps)$ is open.

  By the very same definition, it is clear that $f\big(\CR_f(p,\eps))\subset \CR_f(p,\eps)$. Now, let $y$ be an arbitrary point of $\overline{\CR_f(p,\eps)}$. Since $M$ is compact and $f$ is continuous, there is $\delta>0$ such that $d\big(f(x_1),f(x_2)\big)<\eps$, whenever $d(x_1,x_2)<\delta$. Then, there is $x\in\CR_f(p,\eps)\cap B_\delta(y)$. So there is a finite segment of $\eps$-pseudo orbit $x_0=p, x_1,\ldots,x_n=x$. Since $x\in B_\delta(y)$, we conclude $x_0=p, x_1,\ldots,x_n=x, x_{n+1}:=f(y)$ is a finite segment of $\eps$-pseudo orbit, too. Hence, $f(y)\in\CR_f(p,\eps)$, and we conclude that $f\big(\overline{\CR_f(p,\eps)}\big)\subset \CR_f(p,\eps)$.
\end{proof}

\begin{exercise}
  Let $f\colon M\carr$ be a homeomorphism on the compact metric space $(M,d)$ and $U\subset M$ be a trapping region for $f$. Show that
  \begin{displaymath}
    \CR(f)\subset A_U^+\sqcup A_U^-,
  \end{displaymath}
  where $\sqcup$ denotes the disjoint union and $A_U^+,A_U^-$ is the dual pair of attractor/repellor associated to $U$.
\end{exercise}

Now we can define the notion of \textbf{chain recurrence classes.} If $f\colon M\carr$ is a homeomorphism on the compact metric space $(M,d)$ and $x,y\in\CR(f)$, we say that they are \textbf{chain equivalent} when $y\in\CR_{f}(x)$ and $x\in\CR_{f}(y)$, \ie~$x,y\in\CR_f(x,\eps)\cap\CR_f(y,\eps)$, for every $\eps>0$.

Then we have the following

\begin{exercise}[Chain-recurrence classes]
  If $f$ and $M$ are as above, let us show that \emph{chain equivalence} determines an equivalence relation on $\CR(f)$. Their equivalence classes are called \textbf{chain-recurrence classes.} Show that every chain-recurrent class is compact and $f$-invariant.

  Moreover, two points $x,y\in\CR(f)$ are chain equivalent if and only if either $x,y\in A_U^+$, or $x,y\in A_U^-$, for every trapping region $U$ for $f$.
\end{exercise}

\begin{exercise}
  If $f$ and $M$ are as above, show that the following set is at mots countable:
  \begin{displaymath}
    \left\{(A_U^+,A_U^-) : U\subset M \text{ is a trapping region for  } f\right\}.
  \end{displaymath}
\end{exercise}

Using all previous exercises and results, one can easily prove the so-called ``Fundamental Theorem of Dynamical Systems'' due to Conley~\cite{ConleyIsolInvSets}:

\begin{theorem}[Existence of complete Lyapunov functions]\label{thm:fund-thm-dyn-sys}
  Given a homeomorphism $f\colon M\carr$ on the compact metric space $(M,d)$, there exists a \emph{complete Lyapunov function} for $f$, \ie~a continuous function $L\colon M\to\R$ such that the following conditions hold:
  \begin{enumerate}
    \item $L\big(f(x)\big)\leq L(x)$, for every $x\in M$,
    \item $L\big(f(x)\big)=L(x)$ if and only if $x\in\CR(f)$,
    \item $L(x)=L(y)$, with $x,y\in\CR(f)$ if and only if $x$ and $y$ are chain equivalent, and,
    \item $L\big(\CR(f)\big)$ is a compact totally disconnected set, \ie~it does not contain any open interval.
  \end{enumerate}
\end{theorem}


\section{Dynamics on the real line}
\label{sec:real-line}

\subsection{Affine maps of $\R$}
\label{sec:affine-maps-R}

For each $a,b\in\R$, we can define the affine map $f_{a,b}\colon\R\carr$ by $f_{a,b}(x):=ax+b$, for every $x\in\R$. We are going to describe the orbit structure of $f_{a,b}$ for every choice of $a,b\in\R$.

If $a=1$ and $b=0$, then $f_{1,0}=id_\R$. If $a=1$ and $b>0$, then $f_{1,b}^n(x)\to+\infty$, as $n\to+\infty$, for every $x\in\R$. If $a=1$ and $b<0$, then $f_{1,b}^n(x)\to-\infty$, as $n\to+\infty$, for every $x\in\R$.

If $a\neq 1$ and $b\in\R$ is arbitrary, then $p_{a,b}=-\frac{b}{a-1}$ is the unique fixed point of $f_{a,b}$. Moreover, if $|a|<1$, then $f_{a,b}^n(x)\to p_{a,b}$, as $n\to+\infty$, for every $x\in\R$. If $|a|>1$, then $\abs{f_{a,b}^n(x)}\to+\infty$, as $n\to+\infty$, for every $x\in\R\setminus\{p_{a,b}\}$.

Finally, if $a=-1$ and $b\in\R$ is arbitrary, then we have $f_{-1,b}^2=id_\R$, and $\Fix(f_{-1,b})=\{p_{-1,b}\}$ and $\Per(f_{-1,b})=\R$.

\subsection{The quadratic family}
\label{sec:quadratic-family}

The dynamics of quadratic polynomials in $\R$ is much more interesting than the one of affine maps. Now, for each triple of real numbers $a,b,c\in\R$, with $a\neq 0$, we can define the quadratic map $f_{a,b,c}\colon\R\carr$ by $f_{a,b,c}(x):=ax^{2}+bx+c$, for every $x\in\R$.

First observe that $\Fix(f_{a,b,c})=\emptyset$ if and only if $(b-1)^2-4ac<0$. In such a case, when $a>0$, it holds $f_{a,b,c}^n(x)\to+\infty$, as $n\to+\infty$, for every $x\in\R$. When $a<0$, it holds $f_{a,b,c}^n(x)\to-\infty$, as $n\to+\infty$, for every $x\in\R$.

So, from now on we are going to assume that $(b-1)^2-4ac\geq 0$ (with $a\neq 0$), and hence $\Fix(f_{a,b,c})\neq\emptyset$.

We are going to introduce a fundamental concept in Dynamical Systems, which corresponds to the notion of isomorphism in the category of topological dynamical systems:

\begin{definition}[Conjugacy]
  Given two metric spaces $M,N$ and two continuous maps $f\colon M\carr$ and $g\colon N\carr$, we say that $f$ and $g$ are \textbf{topologically conjugate} when there exists a homeomorphism $h\colon M\to N$ such that
  \begin{displaymath}
    h\circ f = g\circ h.
  \end{displaymath}

  When $M$ and $N$ are differentiable manifolds and $f\colon M\carr$ and $g\colon N\carr$ are $C^r$ maps, with $r\in\N\cup\{\infty\}$, we say that $f$ and $g$ are \textbf{$C^s$-conjugate} when there exists a $C^s$-diffeomorphism $h\colon M\to N$, with $s\in\N\cup\{\infty\}$ and $s\leq r$.
\end{definition}

To simplify the study of the dynamics of quadratic maps, we are going to show that every quadratic map with fixed points is affinely conjugate to a quadratic map of the form $Q_\mu(x):=\mu x(1-x)$, for some $\mu>0$.

In order to do that, first suppose that $a>0$. If we consider the affine map $h\colon \R\ni x\mapsto -x$, we easily see that $h\circ f_{a,b,c}=f_{-a,b,-c}\circ h$. So, from now on we can assume that $a<0$.

Then, the vertex of the parabola defined by $f_{a,b,c}$ is at the point $v=-\frac{b}{2a}$ and since we are assuming that $a<0$ and $(b-1)^2-4ac\geq 0$, we know that $f_{a,b,c}$ has a fixed point $p\in\R$ with $p<v$. So, there is an affine map $h\colon\R\ni x\mapsto \alpha x + \beta$ such that $h(0)=p$ and $h(1/2)=v$.

\begin{problem}
  If $f_{a,b,c}$, $v$, $p$ and $h$ are as above, show that there is $\mu>0$ such that $h\circ f_{a,b,c}=Q_\mu\circ h$.
\end{problem}

So, from now on we shall concentrate on the dynamics of the family of quadratic maps $\{Q_\mu : \mu>0\}$, which is called the \textbf{quadratic family.} For each $\mu>0$, observe that
\begin{displaymath}
  p_\mu:=\frac{\mu-1}{\mu},
\end{displaymath}
is a fixed point of $Q_\mu$, and hence $\Fix(Q_\mu)=\{0,p_\mu\}$.

We have the following result describing the dynamics in the neighborhood of the fixed points:

\begin{proposition}\label{prop:hyperbolic-fixed-points-dim-one}
  Let $I\subset\R$ be an open interval and $f\colon I\carr$ be a $C^1$ map. If $p\in I$ is a fixed point of $f$ and $\abs{f'(p)}\neq 1$, then the following dichotomy holds:
  \begin{itemize}
    \item If $\abs{f'(p)}<1$, then there is $\delta>0$ such that $f^n(x)\to p$, as $n\to+\infty$, for every $x\in (p-\delta,p+\delta)\subset I$;
    \item if $\abs{f'(p)}>1$, then there is $\delta>0$ such that for every $x\in I$ with $0<\abs{x-p}<\delta$, there is $n\in\N$ such that $\abs{f^n(x)-p}\geq\delta$.
  \end{itemize}
\end{proposition}

\begin{proof}
  In the first case, since $\abs{f'(p)}<1$ and $I$ is open, there is $\delta>0$ and $\lambda\in(0,1)$ such that $\abs{f'(x)}<\lambda$, for every $x\in (p-\delta,p+\delta)\subset I$. By the mean value theorem, given any $x\in(p-\delta,p+\delta)$, it holds
  \begin{displaymath}
    \abs{f(x)-p}\leq\lambda\abs{x-p}<\lambda\delta<\delta.
  \end{displaymath}
  So, $f$ restricted to $(p-\delta,p+\delta)$ is a uniform contraction, and by the Banach fixed point theorem, $f^n(x)\to p$, as $n\to+\infty$, for every $x\in(p-\delta,p+\delta)$.

  Now, let us assume that $\abs{f'(p)}>1$. There is $\delta>0$ and $\lambda>1$ such that $\abs{f'(x)}>\lambda$, for every $x\in(p-\delta,p+\delta)$. By the mean value theorem, given any $x\in(p-\delta,p+\delta)$, if $f^i(x)\in(p-\delta,p+\delta)$, for every $i\in\{0,1,\ldots,m\}$, then it holds
  \begin{displaymath}
    \delta>\abs{f^m(x)-p}\geq\lambda^n\abs{x-p}.
  \end{displaymath}
  So,  $m\leq \log(\delta/\abs{x-p})/\log\lambda$ and the proof is completed.
\end{proof}

\begin{definition}
  \label{def:hyperbolic-fixed-point}
  A fixed point $p$ of a $C^1$ map $f\colon I\carr$ is said to be \textbf{hyperbolic} when $\abs{f'(p)}\neq 1$. If $\abs{f'(p)}<1$, then $p$ is called an \textbf{attractor}, and if $\abs{f'(p)}>1$, then $p$ is called a \textbf{repeller}.
\end{definition}

Returning to the quadratic family, we have that $Q_\mu'(x)=-2\mu x + \mu$, and thus, $Q_\mu'(0)=\mu$ and $Q_\mu'(p_\mu)=2-\mu$.

So, if $\mu\in (0,1)$, $p_\mu<0$, $p_\mu$ is a repeller and $0$ is an attractor. In such a case, one can easily check that $Q_\mu^n(x)\to 0$, as $n\to+\infty$, for every $x\in (p_\mu,q_\mu)$, where $q_\mu$ is the unique point of $(p_\mu,+\infty)$ such that $Q_\mu(q_\mu)=p_\mu$; and $Q_\mu^n(x)\to-\infty$, as $n\to+\infty$, for every $x\in\R\setminus[p_\mu,q_\mu]$.

If $\mu=1$, then $p_\mu=0$ and $0$ turns to be a non-hyperbolic fixed point. In such a case, one can easily check that $Q_1^n(x)\to 0$, as $n\to+\infty$, for every $x\in(0,1)$, and $Q_1^n(x)\to-\infty$, as $n\to+\infty$, for every $x\in\R\setminus[0,1]$.

\begin{exercise}
  Show that if $\mu\in(1,3)$, then $p_\mu\in(0,1)$ is a hyperbolic attractor, while $0$ is a repeller, $Q_\mu^n(x)\to p_\mu$, as $n\to+\infty$, for every $x\in(0,1)$; and $Q_\mu^n(x)\to-\infty$, as $n\to+\infty$, for every $x\in\R\setminus[0,1]$.
\end{exercise}

If $\mu=3$, the fixed point $p_\mu$ turns to be a non-hyperbolic attractor, and one can easily check that $Q_3^n(x)\to p_3$, as $n\to+\infty$, for every $x\in(0,1)$; and $Q_3^n(x)\to-\infty$, as $n\to+\infty$, for every $x\in\R\setminus[0,1]$.

If $\mu\in(3,4)$, then both fixed points of $Q_\mu$ are repeller, the interval $[0,1]$ is a totally invariant set for $Q_\mu$ and the dynamics of $Q_\mu$ on $[0,1]$ is much more complicated (see for instance~\cite{DevaneyIntroChaoticDynSys} for more details on this topic).

And now, we shall concentrate on the case $\mu>4$, paying special attention to the case $\mu>\sqrt{5}+2$. When $\mu>4$, the interval $[0,1]$ is not invariant for $Q_\mu$ anymore. In fact, $Q_\mu(1/2)=\mu/4>1$ and hence there are two points $r_\mu, s_\mu$, with  $0<r_\mu<1/2<s_\mu<1$ such that $Q_\mu(r_\mu)=Q_\mu(s_\mu)=1$. Since $Q_\mu^n(x)\to-\infty$, as $n\to+\infty$, for every $x\in\R\setminus[0,1]$, we conclude that $Q_\mu^n(x)\to-\infty$, as $n\to+\infty$, for every $x\in (r_\mu,s_\mu)$. This implies that
\begin{displaymath}
  [0,1]\supsetneq Q_\mu^{-1}([0,1])\supsetneq Q_\mu^{-2}([0,1])\supsetneq\cdots\supsetneq Q_\mu^{-n}([0,1])\supsetneq\cdots,
\end{displaymath}
and hence
\begin{equation}
  \label{eq:Lambda-mu-invariant-set-Q-mu}
  \Lambda_\mu:=\bigcap_{n\in\N} Q_\mu^{-1}\big([0,1]\big),
\end{equation}
is a compact, non-empty totally invariant set for $Q_\mu$, which does not coincide with the interval $[0,1]$.

In fact, $\Lambda_\mu\subset I_\mu^1\sqcup I_\mu^2$, where $I_\mu^1:=[0,r_\mu]$ and $I_\mu^2:=[s_\mu,1]$. So, if we write $\db{k}:=\{1,2,\ldots,k\}$, for any $k\in\N$, we can define the so-called \textbf{itinerary map}  by $h\colon\Lambda_\mu\to\db{2}^{\N_0}$ where
\begin{displaymath}
  {h(x)}_n:=\begin{cases}
    1, &\text{if } Q_\mu^n(x)\in I_\mu^1,\\
    2, &\text{if } Q_\mu^n(x)\in I_\mu^2,
  \end{cases}\quad\forall n\in\N_0.
\end{displaymath}

In general, given any $k\in\N$ with $k\geq 2$, if we endowed the space $\db{k}^{\N_0}$ with the product topology, then by Thychonov theorem we get that $\db{k}^{\N_0}$ is compact. Moreover, its topology is metrizable and compatible with the topology induced by the metric $d\colon\db{k}^{\N_0}\to [0,1]$ given by
\begin{displaymath}
  d\big((i_n),(j_n)\big):=2^{-\min\{n\in\N_0 : i_n\neq j_n\}},\quad\forall (i_n),(j_n)\in\db{k}^{\N_0},
\end{displaymath}
where we are assuming that $\min\emptyset=+\infty$ and hence $d\big((i_n),(j_n)\big)=0$, if $i_n=j_n$, for every $n\in\N_0$.

We need the following

\begin{definition}[Cantor spaces]
  A metrizable topological space $K$ is said to be a \textbf{Cantor space} when it is compact, \textbf{perfect} (\ie~there is no isolated point) and \textbf{totally disconnected} (\ie~if $C\subset K$ is connected and non-empty, then it is a singleton).
\end{definition}

\begin{exercise}
  Show that $\db{k}^{\N_0}$ is a Cantor space, for every $k\in\N$ with $k\geq 2$.
\end{exercise}

With this topology, it is easy to check that the itinerary map $h\colon\Lambda_\mu\to\db{2}^{\N_0}$ that we defined above is, indeed, continuous. Moreover, since $\mu>\sqrt{5}+2$, there is $\lambda >1$ such that $\abs{Q_\mu'(x)}>\lambda$, for every $x\in I_\mu^1\cup I_\mu^2$. So, given any two distinct points $x,y\in\Lambda_\mu$, there is $n\in\N_0$ such that $Q_\mu^n(x)$ and $Q_\mu^n(y)$ belong to different intervals of the form $I_\mu^1$ or $I_\mu^2$. Hence, we conclude that the itinerary map $h\colon\Lambda_\mu\to\db{2}^{\N_0}$ is injective.

Moreover, the map $Q_\mu\big|_{I_\mu^i}\colon I_\mu^i\to [0,1]$ is a diffeomorphisms for $i=1,2$. So, they admit well defined inverse maps that we will denote by $g_i\colon [0,1]\to I_\mu^i$, for $i=1,2$, which are contraction with Lipschitz constants smaller or equal than $\lambda^{-1}$.  Then, given any sequence $(i_n)\in\db{2}^{\N_0}$, we can define the point $x\in\Lambda_\mu$ by the following intersection:
\begin{displaymath}
  \{x\}:=\bigcap_{n\in\N_0} g_{i_0}\circ g_{i_1}\circ\cdots\circ g_{i_n}([0,1]).
\end{displaymath}

One can easily check that $h(x)=(i_n)$, and hence the itinerary map $h\colon\Lambda_\mu\to\db{2}^{\N_0}$ is continuous, injective and surjective, and thus it is a homeomorphism.

On the other hand, for any $k\in\N_0$ larger than $1$, we can define the \textbf{shift map} $\sigma\colon\db{k}^{\N_0}\carr$ by
\begin{displaymath}
  \sigma\big((i_n)\big):=(j_n),\ \text{with } j_n:=i_{n+1}, \quad\forall (i_n)\in\db{k}^{\N_0},\ \forall n\in\N_0.
\end{displaymath}

The main elementary properties of the shift map are given by the following
\begin{exercise}
  If $k\in\N$ with $k\geq 2$, and $\sigma\colon\db{k}^{\N_0}\carr$ denotes the shift map, then shw that the following properties hold:
  \begin{enumerate}
    \item $\sigma$ is continuous and surjective;
    \item the set $\Per(\sigma)$ is dense in $\db{k}^{\N_0}$;
    \item $\sigma$ is \textbf{topologically transitive,} \ie~there is a point $x\in\db{k}^{\N_0}$ such that $\{\sigma^n(x) : n\in\N_0\}$ is dense in $\db{k}^{\N_0}$;
    \item for every $x\in\db{k}^{\N_0}$, the set
          \begin{displaymath}
            \left\{y\in\db{k}^{\N_0} : \exists n\in\N,\ \sigma^n(y)=x\right\}
          \end{displaymath}
          is dense in $\db{k}^{\N_0}$;
    \item for every $U\subset\db{k}^{\N_0}$ open and non-empty, there is $n\in\N$ such that $\sigma^n(U)=\db{k}^{\N_0}$:
  \end{enumerate}
  \end{exercise}


\section{Dynamics of circle maps}
\label{sec:dynamics-circle-maps}

Our preferred model for the circle is the quotient space $\T:=\R/\Z$. The natural quotient map $\pi\colon\R\to\T$ given by $\pi(x):=[x]=x+\Z$, for every $x\in\R$, can be used to endow $\T$ with the quotient topology, which is, by the definition, the finest topology on $\T$ such that $\pi$ is continuous. This topology is metrizable and compatible with the metric $d\colon\T\times\T\to[0,1/2]$ given by
\begin{equation}
  \label{eq:distance-circle}
  d([x],[y]):=\min\{\abs{\tilde{x}-\tilde{y}} : \tilde x\in[x],\ \tilde y\in[y]\}\quad\forall x,y\in\R.
\end{equation}

One can easily check that $\pi\colon\R\to\T$ is a covering map, \ie~for every $[x]\in\T$ there is an open neighborhood $U$ of $[x]$ such that $\pi^{-1}(U)$ is a disjoint union of open intervals of $\R$, each of them homeomorphic to $U$ via $\pi$, and since $\R$ is simply connected, this covering map is the universal covering of $\T$. So, we have the following

\begin{proposition}
  \label{prop:lifting-property-circle}
  Given any continuous map $f\colon\T\carr$, any $x_0\in\T$, any $\tilde x_0\in\R$ with $\pi(\tilde x_0)=x_0$, and any $\tilde y_0\in\R$ such that $\pi(\tilde y_0)=f(x_0)$, there exists a unique continuous map $\tilde f\colon\R\carr$ such that $\pi\circ \tilde f=f\circ\pi$ and $\tilde f(\tilde x_0)=\tilde y_0$. Moreover, if $f$ is a homeomorphism, then $\tilde f$ is a homeomorphism, too.
\end{proposition}

As an easy consequence of the lifting property, we have that if $\tilde f_1,\tilde f_2\colon\R\carr$ are two lifts of the same continuous map $f\colon\T\carr$, then there is $n\in\Z$ such that $\tilde f_1(x)=\tilde f_2(x)+n$, for every $x\in\R$.

On the other hand, we have the following

\begin{proposition}
  \label{prop:degree-circle-maps}
  If $f\colon\T\carr$ is a continuous, then there is an integer $d\in\Z$, called the \textbf{degree of $f$}, such that given any lift $\tilde f\colon\R\carr$ of $f$, it holds $\tilde f(x+1)=\tilde f(x)+d$, for every $x\in\R$. Moreover, the map
  \begin{displaymath}
    \Delta_{\tilde f}(\tilde x):=\tilde f(\tilde x)-d\tilde x, \quad\forall \tilde x\in\R,
  \end{displaymath}
  is continuous and $\Z$-periodic, \ie~$\Delta_{\tilde f}(\tilde x+n)=\Delta_{\tilde f}(\tilde x)$, for every $\tilde x\in\R$ and every $n\in\Z$.
\end{proposition}

\begin{proof}
  Given any $\tilde x\in\R$, since $\pi\big(\tilde f(\tilde x+1)\big)=f\big(\pi(\tilde x+1)\big)=f\big(\pi(\tilde x)\big)=\pi\big(\tilde f(\tilde x)\big)$, we conclude that there is $d\in\Z$ such that $\tilde f(\tilde x+1)=\tilde f(\tilde x)+d$. Since $\tilde f$ is a continuous map, we conclude that the integer $d$ does not depend on the choice of $\tilde x\in\R$. Moreover, by our remark above about different lifts of the same map, we conclude that $d$ does not depend on the choice of the lift $\tilde f$ either, and the first part of the proposition is proved.

  On the other hand, to prove that the map $\Delta_{\tilde f}$ is $\Z$-periodic, it is enough to check that
  \begin{displaymath}
    \Delta_{\tilde f}(\tilde x+1)=\tilde f(\tilde x+1)-d(\tilde x+1)=\tilde f(\tilde x)+d-d\tilde x-d=\Delta_{\tilde f}(\tilde x),
  \end{displaymath}
  for every $\tilde x\in\R$.
\end{proof}

\begin{exercise}
  Given two continuous maps $f,g\colon\T\carr$, we write $d_f$ and $d_g$ for the degrees of $f$ and $g$, respectively, and $d_{f\circ g}$ for the degree of $f\circ g$. Show that
  \begin{displaymath}
    d_{f\circ g}=d_f\cdot d_g.
  \end{displaymath}
\end{exercise}

Since the degree of the identity map $id_\T$ is clearly $1$, by the previous exercise we conclude that if $f\colon\T\carr$ is a homeomorphism, then its degree is either $1$ or $-1$.



\subsection{Dynamics of circle homeomorphisms}
\label{sec:dynam-circle-homeos}

By the last remark of the previous section, we know that the group of homeomorphisms of the circle $\T$, that shall be denoted by $\Homeo{}(\T)$, can be decomposed as the disjoint union of two subsets
\begin{displaymath}
  \Homeo{}(\T)=\Homeo{+}(\T)\sqcup\Homeo{-}(\T),
\end{displaymath}
where we are writing $\Homeo{+}(\T):=\{f\in\Homeo{}(\T) : \text{the degree of } f \text{ is } 1\}$ and $\Homeo{-}(\T):=\{f\in\Homeo{}(\T) : \text{the degree of } f \text{ is } -1\}$.

The elements of $\Homeo{+}(\T)$ shall be called \textbf{orientation-preserving} homeomorphisms, while the elements of $\Homeo{-}(\T)$ shall be called \textbf{orientation-reversing} homeomorphisms. One can easily check that any lift of an element of $\Homeo{+}(\T)$ is a strictly increasing homeomorphism of $\R$, while any lift of an element of $\Homeo{-}(\T)$ is a strictly decreasing homeomorphism of $\R$.

\begin{exercise}[Dynamics of orientation-reversing circle homeomorphisms]
   Given any $f\in\Homeo{-}(\T)$, show that $\Fix(f)$ contains exactly two points and $\Per(f)=\Fix(f^2)$, where $f^2\in\Homeo{+}(\T)$ is an orientation-preserving circle homeomorphism.
\end{exercise}

\begin{example}[Circle rotations]
  \label{exam:circle-rotations}
  Given any $\alpha\in\R$, we can define the translation map $\tilde{R}_\alpha\colon\R\carr$ by $\tilde{R}_\alpha(x):=x+\alpha$, for every $x\in\R$. The map $\tilde{R}_\alpha$ turns to be a lift of an orientation-preserving circle homeomorphism $R_\alpha\colon\T\carr$ that is given by $R_\alpha([x])=[x+\alpha]$, for every $[x]\in\T$ and is called the \textbf{circle rotation} by $\alpha$. The dynamics of $R_\alpha$ is very simple: if $\alpha\in\Q$, and $\alpha=p/q$ with $p\in\Z$ and $q\in\N$ coprime, $R_\alpha^q=id_{\T}$ and hence, $\Per(f)=\Fix(f)=\T$. On the other hand, if $\alpha\in\R\setminus\Q$, then $R_\alpha$ is \textbf{minimal,} \ie~the positive and negative orbit of ever point of $\T$ is dense in $\T$.
\end{example}

\begin{exercise}
  Show that the circle rotation $R_\alpha$ is minimal, for all $\alpha\in\R\setminus\Q$.
\end{exercise}


So, let us start the study of orientation-preserving circle homeomorphisms. We have the following

\begin{exercise}
  A homeomorphism $\tilde f\in\Homeo{}(\R)$ is a lift of an element of $\Homeo{+}(\T)$ if and only if $\tilde f$ is continuous, strictly increasing and $\tilde f(x+1)=\tilde f(x)+1$, for every $x\in\R$. In such a case, $\tilde f$ commutes with every integer translation, \ie~$\tilde f\circ T_m=T_m\circ\tilde f$, for every $m\in\Z$, where $T_m:\R\ni x\mapsto x+m\in\R$, and the function
  \begin{displaymath}
    \Delta_{\tilde f}(x):=\tilde f(x)-x, \quad\forall x\in\R,
  \end{displaymath}
  is continuous and $\Z$-periodic, \ie~$\Delta_{\tilde f}(x+n)=\Delta_{\tilde f}(x)$, for every $x\in\R$ and every $n\in\Z$. In such a case, we can consider the function $\Delta_{\tilde f}$ can be considered as an element of $C^0(\T,\R)$, and is usually called the \textbf{displacement function} of $\tilde f$.
\end{exercise}

From now on, we shall write
\begin{displaymath}
  \Thomeo_+(\T):=\left\{\tilde f\in\Homeo{}(\R) : \tilde f \text{ is a lift of an } f\in\Homeo{+}(\T) \right\}.
\end{displaymath}

\begin{exercise}
  \label{ex:displacement-function-compositio}
  Show that $\Thomeo_+(\T)$ is a subgroup of $\Homeo{}(\R)$. On the other hand, show that given $\tilde f,\tilde g\in\Thomeo_+(\T)$, it holds
  \begin{displaymath}
    \Delta_{\tilde f\circ\tilde g}=\Delta_{\tilde f}\circ g + \Delta_{\tilde g}.
  \end{displaymath}
  Use this fact to show that
  \begin{displaymath}
    \Delta_{\tilde f^n}=\sum_{i=0}^{n-1} \Delta_{\tilde f}\circ f^i, \quad\forall n\in\N.
  \end{displaymath}
\end{exercise}

We have the following elementary but fundamental result about the \textbf{displacement function} of a lift of an orientation-preserving circle homeomorphism:

\begin{proposition}
  \label{prop:displacement-function-iterates}
  It holds
  \begin{displaymath}
    \osc\Delta_{\tilde f}\leq 1, \quad\forall \tilde f\in\Thomeo_+(\T),
  \end{displaymath}
  where the \textbf{oscillation of a function}  $\varphi\in C^0(\T,\R)$ is given by $\osc\varphi:=\max_{x\in\T} \varphi(x)-\min_{y\in\T}\varphi(y)$.
\end{proposition}

\begin{proof}
  Let us suppose that there is $\tilde f\in\Thomeo_+(\T)$ such that $\osc\Delta_{\tilde f}>1$. Then, sine $\Delta_{\tilde f}$ is $\Z$-periodic, there are $x_m,x_M\in\R$ with $0<x_M-x_m<1$ and such that $\Delta_{\tilde f}(x_m)=\min_\T\Delta_{\tilde f}$ and $\Delta_{\tilde f}(x_M)=\max_\T\Delta_{\tilde f}$. Then, we have  that
  \begin{displaymath}
    \tilde f(x_M)-\tilde f(x_m) = \Delta_{\tilde f}(x_M) - \Delta_{\tilde f}(x_m) + x_M - x_m > 1,
  \end{displaymath}
  and hence, by continuity of $\tilde f$, there are two points $x_1,x_2\in(x_m,x_M)$ such that $\tilde f(x_1)=\tilde f(x_2)+1$, and this implies that $f\big(\pi(x_1)\big)=\pi\big(\tilde f(x_1)\big)=\pi\big(\tilde f(x_2)+1\big)=\pi\big(\tilde f(x_2)\big)=f\big(\pi(x_2)\big)$. On the other hand, since $\abs{x_2-x_1}<x_M-x_m<1$, we have $\pi(x_1)\neq\pi(x_2)$, which contradicts the fact that $f$ is injective.
\end{proof}

Proposition~\ref{prop:displacement-function-iterates} lets us define the most important dynamical invariant of orientation-preserving circle homeomorphisms:

\begin{theorem}
  \label{thm:rotation-number-def}
  Given any $\tilde f\in\Thomeo_+(\T)$, there is a number $\rho(\tilde f)\in\R$ such that
  \begin{equation}
    \label{eq:rotation-number-def}
    \rho(\tilde f)=\lim_{n\to+\infty}\frac{\tilde f^n(x)-x}{n}=\lim_{n\to+\infty} \frac{\Delta_{\tilde f^n}(x)}{n}, \quad\forall x\in\R,
  \end{equation}
  where the limit is uniform in $x\in\R$. On the other hand, for every $m,n\in\Z$, it holds
  \begin{displaymath}
    \rho(\tilde f^n\circ T_m)=\rho(\tilde f^n)+m=n\rho(\tilde f)+m,
  \end{displaymath}
  where $T_m\in\Thomeo_+(\T)$ denotes the translation given by $T_m(x)=x+m$, for every $x\in\R$.
\end{theorem}

In order to prove Theorem~\ref{thm:rotation-number-def}, we need the so called \textbf{subadditive lemma}:

\begin{lemma}\label{lem:quasi-morphisms}
  Let ${(a_{n})}_{n\in\N}$ be a sequence of real numbers such that there is a real constant $C>0$ satisfying
  \begin{displaymath}
    \abs{a_{m+n}-a_{m}-a_{n}}\leq C, \quad\forall m,n\in\N.
  \end{displaymath}
  Then, the following limit exists
  \begin{displaymath}
    \rho:=\lim_{n\to+\infty}\frac{a_{n}}{n},
  \end{displaymath}
  and it holds $\abs{a_{n}-n\rho}\leq C$, for every $n\in\N$.
\end{lemma}

\begin{exercise}
  Prove Lemma~\ref{lem:quasi-morphisms}.
\end{exercise}

\begin{proof}[Proof of Theorem~\ref{thm:rotation-number-def}]
  Let us start proving that the limit~\eqref{eq:rotation-number-def} exists for $x=0$. In fact, this is an easy consequence of Proposition~\ref{prop:displacement-function-iterates} and Lemma~\ref{lem:quasi-morphisms}. Indeed, by Proposition~\ref{prop:displacement-function-iterates}, we have that
  \begin{displaymath}
    \begin{split}
      \abs{\Delta_{\tilde f^{m+n}}(0)-\Delta_{\tilde f^m}(0)-\Delta_{\tilde f^n}(0)}&= \abs{\tilde f^{m+n}(0)-\tilde f^m(0)-\tilde f^n(0)} \\
      &= \abs{\Delta_{\tilde f^n}\big(\tilde f^m(0)\big) - \Delta_{\tilde f^n}(0)}\leq\osc\Delta_{\tilde f^n} \leq 1, \\
      \end{split}
    \end{displaymath}
    for every $m,n\in\N$. So, by Lemma~\ref{lem:quasi-morphisms}, there is $\rho(\tilde f)\in\R$ such that $\lim_{n\to+\infty} \frac{\Delta_{\tilde f^n}(0)}{n}=\rho(\tilde f)$.

    Invoking again Proposition~\ref{prop:displacement-function-iterates}, we have that
    \begin{displaymath}
      \abs{\frac{\Delta_{\tilde f^n}(x)}{n}-\frac{\Delta_{\tilde f^n}(0)}{n}}\leq\frac{\osc\Delta_{\tilde f^n}}{n}\leq\frac{1}{n},
    \end{displaymath}
    and hence, the limit is uniform in $x\in\R$.

    Finally, given any $m,n\in\Z$, by Exercise~\ref{ex:displacement-function-compositio}, we have that
    \begin{displaymath}
      \begin{split}
        \rho(\tilde f^n\circ T_m) &= \lim_{k\to+\infty} \frac{(\tilde f^n\circ T_m)^k(x) -x}{k} = \lim_{k\to+\infty}
                                    \frac{\tilde f^{nk}(x)-x + mk}{k} \\
                                  &= m+ \lim_{k\to+\infty} n\frac{\tilde f^{nk}(x)-x}{nk} = n\rho(\tilde f) + m.
      \end{split}
    \end{displaymath}
\end{proof}

We will need the following elementary but fundamental property of the rotation number:

\begin{proposition}
  \label{prop:existence-point-exact-displacement-rotation-number}
  If $\tilde f\in\Thomeo_+(\T)$ and $\rho(\tilde f)$ denotes the rotation number of $\tilde f$, then for each $n\in\N$ there exists $x_n\in\R$ such that
  \begin{displaymath}
    \tilde f^n(x_n)=x_n+n\rho(\tilde f).
  \end{displaymath}
\end{proposition}

\begin{proof}
  Let us fix $n\in\N$. Suppose there is no $x\in\R$ such that $\tilde f^n(x)=x+n\rho(\tilde f)$. Then, let us assume that $\tilde f^n(x)>x+n\rho(\tilde f)$. Since $\Delta_{\tilde f^n}=\tilde f^n-id_{\R}$ is $\Z$-periodic, there exists $\eps>0$ such that $\tilde f^n(x)\geq x+n\rho(\tilde f)+\eps$, for every $x\in\R$. So, by Proposition~\ref{prop:displacement-function-iterates} we have that
  \begin{displaymath}
    \begin{split}
      \rho(\tilde f)&=\lim_{k\to+\infty} \frac{\tilde f^{nk}(x)-x}{nk} =\lim_{k\to+\infty} \frac{\sum_{i=0}^{k-1} \Delta_{\tilde f^n}\big(\tilde f^{ni}(x)\big)}{nk} \\
                    &\geq \lim_{k\to+\infty} \frac{nk\rho(\tilde f)+k\eps}{nk}=\rho(\tilde f)+\frac{\eps}{n},
    \end{split}
  \end{displaymath}
  which is a contradiction.
\end{proof}

We have two important consequences Theorem~\ref{thm:rotation-number-def} and Proposition~\ref{prop:existence-point-exact-displacement-rotation-number}:

\begin{corollary}[Existence of periodic points]
  \label{cor:rational-rotation-number-periodic-point}
  If $f\in\Homeo{+}(\T)$ and $\Per(f)\neq\emptyset$, then $\rho(\tilde f)\in\Q$, for every lift $\tilde f$ of $f$. Conversely, if $\rho(\tilde f)=p/q\in\Q$, with $p\in\Z$ and $q\in\N$ coprime, then there is $\tilde x\in\R$ such that $\tilde f^q(\tilde x)=\tilde x+p$, and every periodic point for $f$ has prime period $q$.
\end{corollary}

\begin{proof}
  If $x\in\T$ is a periodic point for $f$ with prime period $q\in\N$, then for every $\tilde x\in\pi^{-1}(x)$, there is $p\in\Z$ such that $\tilde f^q(\tilde x)=\tilde x+p$. By Proposition~\ref{prop:displacement-function-iterates} we know that $\tilde f$ commutes with every integer translation, and hence, the number $p$ does not depend on the choice of $\tilde x\in\pi^{-1}(x)$. So, if we compute the rotation number of $\tilde f$ using the point $\tilde x$, we get that $\rho(\tilde f)=\lim_{n\to+\infty} \frac{\tilde f^{nq}(\tilde x)-\tilde x}{nq}=\lim_{n\to+\infty} \frac{np}{nq}=p/q$.

  Conversely, if $\rho(\tilde f)=p/q\in\Q$, with $p\in\Z$ and $q\in\N$ coprime, then by Proposition~\ref{prop:existence-point-exact-displacement-rotation-number} there is $\tilde x\in\R$ such that $\tilde f^q(\tilde x)=\tilde x+p$. So, the point $x:=\pi(\tilde x)$ is a periodic point for $f$, where $q$ is a multiple of the prime period of $x$. By the previous argument, if the prime period of $x$ were strictly smaller than $q$, the rotation number of $\tilde f$ would be a rational number with denominator strictly smaller than $q$, which contradicts the fact that $p/q$ is written in irreducible form.
\end{proof}

As an immediate consequence of Proposition~\ref{prop:displacement-function-iterates}, Theorem~\ref{thm:rotation-number-def} and Proposition~\ref{prop:existence-point-exact-displacement-rotation-number}, we have the following fundamental result

\begin{corollary}
  \label{cor:bounded-rotational-deviations-and-continuity-rotation-number}
  Given any $\tilde f\in\Thomeo_+(\T)$, it holds
  \begin{displaymath}
    \abs{\tilde f^n(x)-x-n\rho(\tilde f)}\leq 1, \quad\forall x\in\R,\ \forall n\in\Z.
  \end{displaymath}

  As a consequence of this estimates, we get that the rotation number function $\rho\colon\Thomeo_+(\T)\to\R$ is continuous, when $\Thomeo_+(\T)$ is endowed with the uniform convergence topology.
\end{corollary}

\begin{proof}
  First observe that for every $n\in\N$, there exists $x_n\in\R$ such that $\Delta_{\tilde f^n}(x_n)=n\rho(\tilde f)$. On the other hand, by Proposition~\ref{prop:displacement-function-iterates}, we have that $\osc\Delta_{\tilde f^n}\leq 1$, for every $n$. So, we conclude $\abs{\tilde f^n(x)-x-n\rho(\tilde f)}=\abs{\Delta_{\tilde f^n}(x)-\Delta_{\tilde f^n}(x_n)}\leq\osc\Delta_{\tilde f^n}\leq 1$, for every $n\in\N$. The case, $n<0$ easily follows by applying the previous estimates to the homeomorphism $\tilde f^{-1}$.

  In order to prove that the rotation number function is continuous, let us fix $\tilde f\in\Thomeo_+(\T)$ and $\eps>0$. Let $n$ be a positive integer such that $3/n<\eps$. Then, let us choose a real number $\delta>0$ such that
  \begin{displaymath}
    \abs{\tilde f^j(x)-\tilde g^j(x)}<1, \quad\forall x\in\R,\ \forall j\in\{1,2,\ldots,n\},
  \end{displaymath}
  for every $\tilde g\in\Thomeo_+(\T)$ such that $\abs{\tilde g(x)-\tilde f(x)}<\delta$, for every $x\in\R$. For such a $\tilde g$, we have that
  \begin{displaymath}
    \begin{split}
      \abs{\rho(\tilde f)-\rho(\tilde g)}&\leq \abs{\rho(\tilde f)-\frac{\tilde f^n(x)-x}{n}} + \abs{\frac{\tilde f^n(x)-x}{n}-\frac{\tilde g^n(x)-x}{n}} + \abs{\frac{\tilde g^n(x)-x}{n}-\rho(\tilde g)} \\
                                         &\leq \frac{1}{n} + \frac{1}{n}\abs{\tilde f^n(x)-\tilde g^n(x)} + \frac{1}{n} < \frac{3}{n} < \eps,
    \end{split}
  \end{displaymath}
  and the continuity of the rotation number function is proved.
\end{proof}


We shall analyze now the dynamics of orientation-preserving circle homeomorphisms with irrational rotation number. We have the following fundamental

\begin{theorem}
  \label{thm:irrational-rotation-number-minimality}
  If $f\in\Homeo{+}(\T)$ and $\rho(\tilde f)\in\R\setminus\Q$, for some (and hence every) lift $\tilde f$ of $f$, then the following properties hold:
  \begin{enumerate}
    \item there exists a unique $f$-invariant \textbf{minimal set} $K\subset\T$ for $f$, \ie~$K$ is compact, non-empty, $f(K)=K$, and there is no proper subset of $K$ with these properties;
    \item the set $K$ is either a Cantor set or the whole circle $\T$;
    \item $f$ is a topological extension of the irrational rotation $R_{\rho(\tilde f)}\colon\T\ni x+\Z\mapsto x+\rho(\tilde f)+\Z\in\T$, for every $x\in\R$;
    \item $\Omega(f)=K$ and $\CR(f)=\T$;
    \item $f$ is topologically conjugate to the irrational rotation $R_{\rho(\tilde f)}$ if and only if $K=\T$.
  \end{enumerate}
\end{theorem}

\begin{proof}[Proof of Theorem~\ref{thm:irrational-rotation-number-minimality}]
  By Exercise~\ref{ex:inclusions-fix-per-rec-limit-sets-non-wandering} we know that there is a minimal set $K\subset\T$ for $f$. Reasoning by contradiction, let us assume that there are another one $K'\subset\T$ with $K\neq K'$. Since $K\cap K'$ is a compact $f$-invariant set that is contained in $K$ and $K'$, we conclude that $K\cap K'=\emptyset$. So, there is $\eps>0$ such that $d(x,y)>\eps$, for every $x\in K$ and every $y\in K'$. On the other hand, let $y_0$ be an arbitrary point of $K'$ and let us write $I_0$ for the connected component of $\T\setminus K$ that contains $y_0$. Since $K$ is $f$-invariant, $f^n(I_0)$ is a connected component of $\T\setminus K$, for every $n\in\Z$; and since $f$ is periodic point free, $f^m(I_0)\cap f^n(I_0)=\emptyset$, for every $m,n\in\Z$ with $m\neq n$. If we write $\abs{f^m(I_0)}$ for the length of the interval $f^m(I_0)$, we have that $\sum_{m\in\Z} \abs{f^m(I_0)}\leq 1$, and hence, $\lim_{m\to\pm\infty} \abs{f^m(I_0)}=0$. By convergence of the series, there is $m_0\in\Z$ such that $\abs{f^{m_0}(I_0)}<\eps$ and $f^{m_0}(y_0)\in K'\cap f^{m_0}(I_0)$, contradicting the fact that $d(x,y)>\eps$, for every $x\in K$ and every $y\in K'$. So, there is a unique minimal set $K\subset\T$ for $f$.

  To prove the second point, let us suppose that $K\neq\T$. Since $K$ is a minimal set, one can easily check that any isolated point of $K$ should be a periodic point for $f$, and since $f$ is periodic point free, we conclude that $K$ has no isolated points. On the other hand, if $J\subset K$ were a non-trivial connected component of $K$, since $K$ is a minimal set, we would have that there exists $n\in\N$ such that $f^n(J)=J$, and hence, $f^n$ would have a fixed point in $J$, contradicting the fact that $f$ is periodic point free. So, $K$ has no non-trivial connected components, and hence, it is a Cantor set.

  In order to prove that $f$ is a topological extension of the irrational rotation $R_{\rho(\tilde f)}$, we first need the following

  \begin{theorem}[Gottschalk-Hedlund theorem]
    \label{thm:Gottschalk-Hedlund}
    Let $(X,d)$ be a compact metric space, $f\colon X\carr$ be a minimal homeomorphism, and $\varphi\in C^0(X,\R)$. Then, there is $u\in C^0(X,\R)$ such that $\varphi=u\circ f-u$ if and only if there is $x_0\in X$ such that the sequence of Birkhoff sums satisfies
    \begin{displaymath}
      \sup_{n\in\N} \abs{\sum_{i=0}^{n-1} \varphi\big(f^i(x_0)\big)}<+\infty.
    \end{displaymath}
  \end{theorem}

  \begin{exercise}
    Prove Theorem~\ref{thm:Gottschalk-Hedlund}.
  \end{exercise}

  Now, let us consider the function $\varphi\in C^0(\T,\R)$ given by $\varphi(x):=\Delta_{\tilde f}(\tilde x)-\rho(\tilde f)$, where $\tilde x$ is any lift of $x\in\T$. Combining Corollary~\ref{cor:bounded-rotational-deviations-and-continuity-rotation-number} and Theorem~\ref{thm:Gottschalk-Hedlund}, we conclude there is $u\in C^0(K,\R)$ such that
  \begin{displaymath}
    \varphi(x)=u\big(f(x)\big)-u(x), \quad\forall x\in K.
  \end{displaymath}

  We use function $u$ to define a continuous functions $\tilde h\colon\tilde{K}\to\R$ and $h\colon K\to\T$ given by $\tilde h(\tilde x):=\tilde x+u(x)$ and $h(x):=x+\pi\big(u(x)\big)$, for every $\tilde x\in\tilde{K}:=\pi^{-1}(K)\subset\R$ and $x=\pi(\tilde x)$. One can easily check $h\circ f\big|_K=R_{\rho(\tilde f)}\circ h$.

  In the case $K=\T$, one can already finish the proof of the theorem:

  \begin{exercise}
    Assuming $K=\T$, show that $h\colon\T\carr$ is injective and surjective, and hence, it is a topological conjugacy between $f$ and $R_{\rho(\tilde f)}$.
  \end{exercise}

  It remains to show that $h$ can be extended to the whole circle $\T$ when $K\neq\T$. In order to do that, let $I$ be an arbitrary connected component of $\T\setminus K$, let $x_0,x_1\in K$ be extreme points of $I$. We are going to show that $h(x_0)=h(x_1)$, and hence, we can extend $h$ to $I$ as the constant function on $I$ with value $h(x_0)=h(x_1)$. To do that, let us suppose that $h(x_0)\neq h(x_1)$ and write $\eps:=d\big(h(x_0),h(x_1)\big)>0$. Since $h\colon K\to\T$ is continuous and $K$ is compact, $h$ is uniformly continuous and hence there is $\delta>0$ such that $d\big(h(x),h(y)\big)<\eps/2$, for every $x,y\in K$ with $d(x,y)<\delta$. Observe that since $h\circ f\big|_K=R_{\rho(\tilde f)}\circ h$, we have that $d\big(h\big(f^n(x_0)\big),h\big(f^n(x_1)\big)\big)=d\big(R_{\rho(\tilde f)}^n\big(h(x_0)\big),R_{\rho(\tilde f)}^n\big(h(x_1)\big)\big)=d\big(h(x_0),h(x_1)\big)=\eps$, for every $n\in\Z$, where $f^n(x_0)$ and $f^n(x_1)$ are the extreme points of the connected component $f^n(I)$ of $\T\setminus K$. But, by our previous argument about the dynamics of connected components of $\T\setminus K$, we know that $\abs{f^n(I)}\to 0$, as $n\to\pm\infty$. So, there is $n\in\Z$ such that $d\big(f^n(x_0),f^n(x_1)\big)<\delta$, contradicting  the property of uniform continuity of $h$.
\end{proof}



\subsection{$C^r$-generic theory of $\Diff{}r(\T)$}
\label{sec:Cr-generic-theory-circle}

In this section we start the study of the so-called $C^r$-generic theory of circle diffeomorphisms, where $r\in\N_0\cup\{\infty\}$ denotes a fixed regularity. We will denote by $\Diff{}r(\T)$ the group of $C^r$-diffeomorphisms of the circle $\T$, and $\Diff+r(\T)$ and $\Diff-r(\T)$ for the connected components of orientation-preserving and orientation-reversing circle diffeomorphisms, respectively.

We endow $\Diff+r(\T)$ with the $C^r$-topology, which is the topology induced by the $C^r$-metric that is inductively defined as follows: for $r=0$, we consider the distance $d_{C^0}\colon\Homeo{}(\T)\times\Homeo{}(\T)\to[0,1/2]$ given by
\begin{displaymath}
  d_{C^0}(f,g):=\max_{x\in\T} d\big(f(x),g(x)\big)+d\big(f^{-1}(x),g^{-1}(x)\big), \quad\forall f,g\in\Homeo{}(\T),
\end{displaymath}
where $d\colon\T\times\T\to[0,1/2]$ is distance function given by~\eqref{eq:distance-circle}. For $r\geq 1$, we endow $\Diff+r(\T)$ with the $C^r$-metric $d_{C^r}\colon\Diff+r(\T)\times\Diff+r(\T)\to[0,+\infty)$ given by
\begin{equation}
  \label{eq:Cr-metric-circle}
  d_{C^r}(f,g):=d_{C^{r-1}}(f,g)+\max_{x\in\T} \abs{\Delta_{\tilde f}^{(r)}(x)-\Delta_{\tilde g}^{(r)}(x)}, \quad\forall f,g\in\Diff+r(\T),
\end{equation}
where $\phi^{(r)}$ denotes the $r$-th derivative of a function $\phi\in C^r(\T,\R)$, and $\tilde f,\tilde g\in\Thomeo_+(\T)$ are any lifts of $f$ and $g$, respectively. One can analogously define the $C^r$-metric on $\Diff-r(\T)$.

We need the following definition that extends the previous one of periodic fixed points (Definition~\ref{def:hyperbolic-fixed-point}):
\begin{definition}
  Let $f\in\Diff{}1(\T)$ and $p\in\T$ be a periodic point for $f$ with prime period $n\in\N$. We say that $p$ is a \textbf{simple periodic point} for $f$ if $Df^n(p)\neq 1$. Moreover, we say that $p$ is \textbf{hyperbolic} if $\abs{Df^n(p)}\neq 1$. A hyperbolic periodic point is said to be \textbf{attracting} if $\abs{Df^n(p)}<1$, and \textbf{repelling} if $\abs{Df^n(p)}>1$.
\end{definition}

We can now define a fundamental family of circle diffeomorphisms:
\begin{definition}[Morse-Smale circle diffeomorphisms]
  A diffeomorphism $f\in\Diff{}1(\T)$ is said to be a \textbf{Morse-Smale} when $\Per(f)\neq\emptyset$ and every periodic point for $f$ is hyperbolic.
\end{definition}

We can now state the main result of this section, which is the following

\begin{theorem}[Morse-Smale circle diffeomorphisms]
  \label{thm:Morse-Smale-circle-diffeomorphisms}
  Given any $r\in\N\cup\{\infty\}$, if we write $MS^r$ for the set of orientation-preserving Morse-Smale circle diffeomorphisms, then $MS^r$ is $C^1$-open (\ie~it is an open set when $\Diff{+}r(\T)$ is endowed with the $C^1$-distance) and $C^r$-dense in $\Diff{+}r(\T)$.
  The set of Morse-Smale circle diffeomorphisms is open and dense in $\Diff{}1(\T)$.
\end{theorem}

We shall need several partial results in order to prove Theorem~\ref{thm:Morse-Smale-circle-diffeomorphisms}. We start with the following one, which states the existence of the so-called \textbf{continuation of simple periodic points}:

\begin{exercise}
  \label{ex:continuation-simple-periodic-points}
  Let $f\colon\T\carr$ be a $C^1$-map and $p\in\T$ be a simple periodic point for $f$ with prime period $n\in\N$. Show that there is an open neighborhood $\mc{U}$ of $f$ in $C^1(\T,\T)$, a real number $\eps>0$, and a continuous map $\mc{U}\ni g\mapsto p_g\in B_\eps(p)\subset\T$ such that $p_f=p$ and $p_g$ is the only periodic point for $g$ in the ball $B_\eps(p)$ and it has with prime period $n$, for every $g\in\mc{U}$. Moreover, if $p$ is hyperbolic, then $p_g$ is hyperbolic with the same type as $p$, for every $g\in\mc{U}$.
\end{exercise}

Next, we need to show that any orientation-preserving $C^r$-diffeomorphism with irrational rotation number can be $C^r$-approximated by diffeomorphisms with rational rotation number:

\begin{theorem}[$C^\infty$-closing lemma on the circle]
  \label{thm:closing-lemma-circle}
  Let $f\in\Homeo+(\T)$ be an arbitrary orientation-preserving circle homeomorphism, and $\tilde f\in\Thomeo_+(\T)$ be any lift of $f$. For each $t\in\R$, consider the homeomorphism $f_t\in\Homeo+(\T)$ given by $\tilde f_t(x):=\tilde f(x)+t$, for every $x\in\R$. Then, the function $\R\ni t\mapsto \rho(\tilde f_t)$ is continuous and if $\rho(\tilde f)\in\R\setminus\Q$, then for every $\eps>0$ there is $t_-\in (-\eps,0)$ and $t_+\in (0,\eps)$ such that $\rho(\tilde f_{t_-})<\rho(\tilde f)<\rho(\tilde f_{t_+})$ and $\rho(\tilde f_{t_-}),\rho(\tilde f_{t_+})\in\Q$.
\end{theorem}

\begin{proof}[Proof of Theorem~\ref{thm:closing-lemma-circle}]
  By Corollary~\ref{cor:bounded-rotational-deviations-and-continuity-rotation-number} we know the function $\R\ni t\mapsto \rho(\tilde f_t)$ is continuous. On the other hand, one can easily check by induction on $n\in\N$ that
  \begin{equation}
    \label{eq:monotonicity-rotation-perturbations}
    \tilde f_{-t}^n(\tilde x)\leq \tilde f^n(\tilde x)\leq \tilde f_t^n(\tilde x), \quad\forall \tilde x\in\R,\ \forall t\geq 0.
  \end{equation}
  In particular, this implies the function $t\mapsto \rho(\tilde f_t)$ is non-decreasing.

  Now, let us assume that $\rho(0)=\rho(\tilde f)\in\R\setminus\Q$. By Theorem~\ref{thm:irrational-rotation-number-minimality}, there is a unique minimal set $K\subset\T$ for $f$, which is either a Cantor set or the whole circle $\T$. In any case, there exists $x\in K$ such that $(\tilde x,\tilde x+\eps)\cap\tilde K\neq\emptyset$ and $(\tilde x-\eps,\tilde x)\cap\tilde K\neq\emptyset$, for every $\eps>0$, where $\tilde K:=\pi^{-1}(K)$.

  Now, let us show that for every $\eps>0$ there is $t_+\in (0,\eps)$ such that $\rho(\tilde f_{t_+})>\rho(\tilde f)$. Since $K$ is a minimal set for $f$ and $x$ is accumulated by both sides in $K$, we know that there is $n\in\N$ and $p\in\Z$ such that
  \begin{equation}
    \label{eq:closing-point-x-lift}
    \tilde x+p-\eps<\tilde f^n(\tilde x)<\tilde x + p.
  \end{equation}

  Combining estimates~\eqref{eq:monotonicity-rotation-perturbations} and~\eqref{eq:closing-point-x-lift}, we conclude that
  \begin{displaymath}
    \tilde f_\eps^n(\tilde x) \geq \tilde f^n(\tilde x) + \eps > \tilde x + p,
  \end{displaymath}
  So, by continuity of $\tilde f_t$ on the parameter $t$, we conclude there is $t_+\in(0,\eps)$ such that $\tilde f_{t_+}^n(\tilde x)=\tilde x+p$, and this implies $\rho(\tilde f_{t_+})>\rho(\tilde f)$. The other case showing the existence of $t_-\in (-\eps,0)$ with $\rho(\tilde f_{t_-})\in\Q$ and $\rho(\tilde f_{t_-})<\rho(\tilde f)$ is completely analogous.
\end{proof}

Then we need to discuss some notions of \textbf{transversality}: given a smooth manifold $M$, we say that two submanifolds $N_1,N_2\subset M$ are \textbf{transverse at the point $x\in N_1\cap N_2$} if it holds
\begin{displaymath}
  T_xM=T_xN_1 + T_xN_2.
\end{displaymath}
We say that they are just \textbf{transverse} if they are transverse at every point of $N_1\cap N_2$. Similarly, given a smooth map $f\colon M\to N$ between smooth manifolds, and a submanifold $P\subset N$, we say that $f$ is \textbf{transverse to $P$ at the point $x\in M$} when $f(x)\in P$ and
\begin{displaymath}
  T_{f(x)}N=T_{f(x)}P + Df_x(T_xM).
\end{displaymath}
We say that $f$ is \textbf{transverse to $P$} if it is transverse to $P$ at every point $x\in M$ such that $f(x)\in P$. Observe that, by definition, $f$ is transverse to $P$ whenever $f(M)\cap P=\emptyset$.

Given smooth manifolds $\Lambda$, $M$ and $N$, a closed submanifold $S\subset N$, and a $C^r$ map $F\colon \Lambda\times M\to N$, if we consider the manifold $\Lambda$ as a parameter space, we can define the set $T_{\Lambda,S}$ by
\begin{equation}
  \label{eq:T-lambda-S-transverse-param}
  T_{\Lambda,S}:=\left\{\lambda\in\Lambda  : F(\lambda,\cdot )\colon M\to N \text{ is transverse to } S\right\}.
\end{equation}
We have the following fundamental result about \emph{parametric transversality:}

\begin{theorem}
  \label{thm:transversality-parametric}
  If $\Lambda$, $M$, $N$ and $S$ are as above and $F\colon\Lambda\times M\to N$ is a $C^r$ map with $r>\max\{0,\dim N+\dim A-\dim M\}$, and $F$ is transverse to $S$, then the set $T_{\Lambda,S}$ given by~\eqref{eq:T-lambda-S-transverse-param} is residual, and hence, dense in $\Lambda$.
\end{theorem}

\begin{proof}
  See~\cite[Theorem 2.7]{HirschDiffTopology}.
\end{proof}

We shall also need the following

\begin{exercise}
  \label{exer:simple-periodic-points-perturbation}
  Let $f\colon\T\carr$ be a $C^r$ map (with $r\geq 1$) $p\in\T$ a periodic point of prime period $n$ for $f$. Show that given any $\eps>0$, there is a $C^\infty$ map $g\colon\T\carr$ to with $d_{C^r}(f,g)<\eps$, where $d_{C^r}$ denotes the $C^r$-distance given by~\eqref{eq:Cr-metric-circle}, and such that $p$ is a simple periodic point for $g$ with prime period $n$.
\end{exercise}

\begin{proof}[Proof of Theorem~\ref{thm:Morse-Smale-circle-diffeomorphisms}]
  Let us fix $r\in\N\cup\{\infty\}$ and let us write $MS^r$ for the set of orientation-preserving Morse-Smale circle diffeomorphisms. First, let us show that $MS^r$ is $C^1$-open. Let $f\in MS^r$ and let $\Per(f)=\{p_1,p_2,\ldots,p_k\}$ be the set of periodic points for $f$. We know that all of them have prime period $q\in\N$. Since every periodic point for $f$ is hyperbolic, by Exercise~\ref{ex:continuation-simple-periodic-points}, for each $i\in\{1,\ldots,k\}$ there is an open neighborhood $\mc{U}_i$ of $f$ in $C^1(\T,\T)$, a real number $\eps_i>0$, and a continuous map $\mc{U}_i\ni g\mapsto p_{i,g}\in B_{\eps_i}(p_i)\subset\T$ such that $p_{i,f}=p_i$ and $p_{i,g}$ is the only periodic point for $g$ in the ball $B_{\eps_i}(p_i)$ and it has with prime period equal to the prime period of $p_i$, for every $g\in\mc{U}_i$. Moreover, every $p_{i,g}$ is hyperbolic.

  On the other hand, there is $\delta>0$ such that
  \begin{displaymath}
    d\big(f^q(x),x)>\delta, \quad\forall x\in\T\setminus\bigcup_{i=1}^k B_{\eps_i}(p_i),
  \end{displaymath}
  where $q$ denotes the prime period of every periodic point of $f$.

  Then, we define the set $\mc{U}$ as the intersection of the neighborhoods $\mc{U}_i$ with the set of maps $g\in\Diff{+}1(\T)$ such that $d(g^q(x),x)>\delta/2$, for every $x\in\T\setminus\bigcup_{i=1}^k B_{\eps_i}(p_i)$. It is easy to check that $\mc{U}$ is a $C^1$-neighborhood of $f$ that is contained in $MS^1$. In particular, $MS^r$ is open when it is endowed with the $C^1$-topology.

  In order to see that $MS^r$ is $C^r$-dense in $\Diff{+}r(\T)$, let us assume $r$ is finite, fix an arbitrary $f\in\Diff{+}r(\T)$ and let $\tilde f\in\Thomeo_+(\T)$ be any lift of $f$. If $\rho(\tilde f)\in\Q$, then by Corollary~\ref{cor:rational-rotation-number-periodic-point} there is a periodic point for $f$. If $\rho(\tilde f)\in\R\setminus\Q$, then by Theorem~\ref{thm:closing-lemma-circle}, for every $\eps>0$ there is $g\in\Diff{+}r(\T)$ with $d_{C^r}(f,g)<\eps$ and such that $\rho(\tilde g)\in\Q$, where $\tilde g$ is any lift of $g$. In any case, we can assume that $f$ has a periodic point $p\in\T$ with prime period $n\in\N$. By Exercise~\ref{exer:simple-periodic-points-perturbation}, for every $\eps>0$ there is $g\in\Diff{+}r(\T)$ with $d_{C^r}(f,g)<\eps$ and such that $p$ is a simple periodic point for $g$ with the very same prime period $n$. Since $p$ is simple, by Exercise~\ref{ex:continuation-simple-periodic-points} we know that the continuation of $p$ is well defined in a $C^1$-neighborhood of $g$. In particular, the rotation number is constant in a $C^1$-neighborhood of $g$. Hence, if $\tilde g\colon\R\carr$ is a lift of $g$, then there is $k\in\Z$ such that $\rho(\tilde g)=k/n$. On the other hand, for each $t\in\R$, consider the diffeomorphism $\tilde g_t:=g+t$. By induction in $n$, one can easily check that the function $G_n\colon\R\times\T\to\T\times\T$ given by
  \begin{displaymath}
    G_n(t,x):=\big(x,\pi\circ\tilde g_t^n(\tilde x)\big), \quad \forall (t,x)\in\R\times\T, \ \forall x\in\pi^{-1}(x),
  \end{displaymath}
  is transverse to the diagonal $D:=\{(y,y)\in\T\times\T : y\in\T\}$; and by our hypothesis about $\rho(\tilde g)$, we know that the image of $G_n$ intersects $D$ at the point $(0,p)$. So, by Theorem~\ref{thm:transversality-parametric}, for every $\eps>0$ there is $t\in(-\eps,\eps)$ such that $G_n(t,x)\in D$, for some $x\in\T$ and the map $G_n(t,\dot)\colon\T\carr$ is transverse to $D$ at every point. In particular,  we get that the $g_t:=R_t\circ g\in MS^r$, and the density, for $r$ finite, is proved. The case $r=\infty$ easily follows by the density of $C^\infty$ maps in $C^r$ maps, for every finite $r$.
\end{proof}

\subsection{Dynamics of circle endomorphisms with $\abs{d}\geq 2$}
\label{sec:dynam-circle-endomorphisms}

Given any $d\in\Z$ and any $r\in\N_0\cup\{\infty\}$, we shall write
\begin{equation}
  \label{eq:degree-d-cirle-maps}
  C_d^r(\T):=\left\{f\colon\T\carr : f \text{ is } C^r,\ \deg f =r\right\}.
\end{equation}

For each $d\in\Z$, we consider the unique group homomorphism $E_d\colon\T\carr$ which is contained in $C_d^\infty(\T)$. This map $E_d$ is given by $E_d(x):=d x$, for every $x\in\T$.

\begin{exercise}[Dynamics of $E_d$]
  Given $d\in\Z$ with $\abs{d}\geq 2$, show the following statements about the dynamics of $E_d$:
  \begin{enumerate}[(i)]
    \item $\Per(E_d)$ is dense in $\T$;
    \item given any non-empty open set $U\subset\T$, there is $n\in\N$ such that $E_d^n(U)=\T$;
    \item given any pair of non-empty open sets $U,V\subset\T$, there is $n\in\N$ such that $E_d^n(U)\cap V\neq\emptyset$;
    \item use the previous points to conclude that there is $x\in\T$ such that $\omega_{E_d}(x)=\T$;
    \item for every $y\in\T$, it holds
          \begin{displaymath}
            \overline{\bigcup_{n\in\N} E_d^{-n}(y)}=\T.
          \end{displaymath}
  \end{enumerate}
\end{exercise}

We have the following fundamental result about the dynamics of circle endomorphisms with degree $\abs{d}\geq 2$:

\begin{theorem}[$E_d$ as topological factor]
  \label{thm:Ed-topological-factor}
  If $d\in\Z$ with $\abs{d}\geq 2$, then every $f\in C_d^0(\T)$ is a topological extension of $E_d\colon\T\carr$, \ie~there is a continuous surjective map $h\colon\T\carr$ such that $h\circ f=E_d\circ h$.
\end{theorem}

\begin{proof}[Proof of Theorem~\ref{thm:Ed-topological-factor}]
  Let $f\in C_d^0(\T)$ be arbitrary. We shall construct a map $h\in C^0_1(\T)$ such that $h\circ f=E_d\circ h$. To do this, let $\tilde f\colon\R\carr$ be a lift of $f$ and notice that $\tilde f(\tilde x)=d\tilde x+\Delta(\tilde x)$, for all $\tilde x\in\R$ and where $\Delta\in C^0(\T,\R)$. Let $\tilde E_d(x):=dx$, for all $x\in\R$, be our chosen lift of $E_d$. Then, we shall look for a lift $\tilde h\colon\R\carr$ with $\tilde h=id+u$ and $u\in C^0(\T,\R)$ satisfying $\tilde h\circ\tilde f=\tilde E_d\circ\tilde h$ and noticing the following functional equation must hold:
  \begin{displaymath}
    d\tilde x + u\big(f(x)\big) + \Delta(x) = d\tilde x + du(x), \quad\forall \tilde x\in\R,
  \end{displaymath}
  and $x=\pi(\tilde x)$, so we get
  \begin{equation}
    \label{eq:functional-equation-u-semi-conjugacy-Ed}
    u(x)= \frac{1}{d} u\big(f(x)\big) + \frac{1}{d} \Delta(x), \quad\forall x\in\T.
  \end{equation}

  Motivated by equation~\eqref{eq:functional-equation-u-semi-conjugacy-Ed}, we define the operator $\mc{L}\colon C^0(\T,\R)\to C^0(\T,\R)$ given by
  \begin{displaymath}
    \mc{L}(u)(x):= \frac{1}{d} u\big(f(x)\big) + \frac{1}{d} \Delta(x), \quad\forall u\in C^0(\T,\R),\ \forall x\in\T.
  \end{displaymath}
  If we endow the space $C^0(\T,\R)$ with the uniform norm $\norm{\phi}:=\max_{x\in\T} \abs{\phi(x)}$, for every $\phi\in C^0(\T,\R)$, then one can easily check that $\mc{L}$ is a contraction, since for every $u,v\in C^0(\T,\R)$ we have
  \begin{displaymath}
    \norm{\mc{L}(u)-\mc{L}(v)} = \norm{\frac{1}{d} u\circ f - \frac{1}{d} v\circ f} \leq \frac{1}{\abs{d}} \norm{u-v}.
  \end{displaymath}

  Hence, by the Banach fixed point theorem, there is a unique $u\in C^0(\T,\R)$ such that $\mc{L}(u)=u$, and this implies that the lift $\tilde h\colon\R\carr$ given by $\tilde h(\tilde x):=\tilde x + u(x)$, for every $\tilde x\in\R$ and $x=\pi(\tilde x)$, satisfies $\tilde h\circ\tilde f=\tilde E_d\circ\tilde h$.
\end{proof}

The map $E_d$, with $\abs{d}\geq 2$ is an example of an \textbf{expanding map:}

\begin{definition}[Expanding circle maps]
  \label{def:expanding-circle-maps}
  A map $f\in C_d^1(\T)$ is said to be \textbf{expanding} when there is $C>0$ and $\lambda>1$ such that
  \begin{displaymath}
    \abs{Df^n(x)}\geq C\lambda^n, \quad\forall x\in\T,\ \forall n\in\N.
  \end{displaymath}
\end{definition}

\begin{exercise}[Degree of expanding maps]
  Let $f\in C_d^1(\T)$ be an expanding map. Show that $\abs{d}\geq 2$.
\end{exercise}

\begin{exercise}
  \label{exer:expanding-maps-equiv-definition}
  Let $f\in C_d^1(\T)$ be an arbitrary map with $\abs{d}\geq 2$. Show that $f$ is an expanding map if and only if there is $n\in\N$ such that
  \begin{displaymath}
    \abs{Df^n(x)}>1, \quad\forall x\in\T.
  \end{displaymath}
\end{exercise}

We have the following fundamental result about the dynamics of expanding circle maps:
\begin{theorem}[Topological conjugacy of expanding circle maps]
  \label{thm:expanding-maps-topological-conjugacy}
  If $f\in C^1_d(\T)$ is an expanding map, then there is a homeomorphism $h\in\Homeo+(\T)$ such that $h\circ f=E_d\circ h$.
\end{theorem}

\begin{proof}
  In Theorem~\ref{thm:Ed-topological-factor} we have shown that there is $h\in C^0_1(\T)$ such that $h\circ f=E_d\circ h$. We are going to show that $h$ is injective, and hence, it is a homeomorphism. To do that, let us assume that there are $x_1,x_2\in\T$ with $x_1\neq x_2$ and $h(x_1)=h(x_2)$. In fact, in the proof of Theorem~\ref{thm:Ed-topological-factor} we have shown that given a lift $\tilde f\colon\R\carr$, there is a lift $\tilde h\colon\R\carr$ of $h$ such that $\tilde h\circ\tilde f=\tilde E_d\circ\tilde h$.
  Let us choose points $\tilde x_i\in\pi^{-1}(x_i)\cap [0,1]$, for $i\in\{1,2\}$. Since $h(x_1)=h(x_2)$, we have $m:=\tilde h(\tilde x_2)-\tilde h(\tilde x_1)\in\Z$. Then, $\tilde x_1+m\neq\tilde x_2$ and $\tilde h(\tilde x_1+m)=\tilde h(x_1)+m=\tilde h(\tilde x_2)$. So, we have
  \begin{displaymath}
    \tilde h\big(\tilde f^n(\tilde x_1+m)\big) = \tilde E_d^n\big(\tilde h(\tilde x_1+m)\big) = \tilde E_d^n\big(\tilde h(\tilde x_2)\big) = \tilde h\big(\tilde f^n(\tilde x_2)\big), \quad\forall n\in\N.
  \end{displaymath}

  Since $f$ is an expanding map, there is $C>0$ and $\lambda>1$ such that $\abs{Df^n(x)}=\abs{D\tilde f^n\big(\pi(x)\big)}\geq C\lambda^n$, for every $x\in\T$ and $n\in\N$. This implies,
  \begin{displaymath}
    \abs{\tilde f^n(\tilde x_1+m)-\tilde f^n(\tilde x_2)}\geq C\lambda^n \abs{\tilde x_1+m-\tilde x_2}\to+\infty, \quad\forall n\in\N,
  \end{displaymath}
  and as $n\to+\infty$. On the other hand, since $\tilde h$ is a lift of a degree one map of the circle, there is $L>1$ such that $\abs{\tilde h(\tilde x)-\tilde h(\tilde y)}\geq L^{-1}\abs{\tilde x-\tilde y}-L$, for every $\tilde x,\tilde y\in\R$. Hence, we get
  \begin{displaymath}
    \begin{split}
      0=\abs{\tilde h\big(\tilde f^n(\tilde x_1+m)\big)-\tilde h\big(\tilde f^n(\tilde x_2)\big)} &\geq L^{-1}\abs{\tilde f^n(\tilde x_1+m)-\tilde f^n(\tilde x_2)} - L\\
      &\geq L^{-1}C\lambda^n \abs{\tilde x_1+m-\tilde x_2} - L \to +\infty, \quad\text{as } n\to+\infty,
    \end{split}
  \end{displaymath}
  which is clearly a contradiction. Hence, $h$ is injective and this concludes the proof.
\end{proof}

If we endow the space $C^r_d(\T)$ with the uniform $C^r$-topology (see, for instance, \cite[Chapter 2]{HirschDiffTopology} for details), as a straightforward consequence of Theorem~\ref{thm:expanding-maps-topological-conjugacy} we get the following

\begin{corollary}[Structural stability of circle expanding maps]
  For every $d\in\Z$ with $\abs{d}\geq 2$, the set
  \begin{displaymath}
    \mc{E}_d:=\left\{f\in C^1_d(\T) : f \text{ is an expanding map}\right\}
  \end{displaymath}
  is $C^1$-open in $C_d^1(\T)$. Thus, every expanding map is $C^1$-structurally stable, \ie~for every $f\in\mc{E}_d$, there is a $C^1$-neighborhood $\mc{U}$ of $f$ in $C_d^1(\T)$ such that every $g\in\mc{U}$ is topologically conjugate to $f$.
\end{corollary}

\begin{proof}
  Let $f\in\mc{E}_d$ be an arbitrary expanding map. So, there is $C>0$ and $\lambda>1$ such that $\abs{Df^n(x)}\geq C\lambda^n$, for every $x\in\T$ and $n\in\N$. In particular, there is $n_0\in\N$ such that $\abs{Df^{n_0}(x)}>1$, for every $x\in\T$. By continuity of the derivative, there is a $C^1$-neighborhood $\mc{U}$ of $f$ in $C_d^1(\T)$ such that $\abs{Dg^{n_0}(x)}>1$, for every $g\in\mc{U}$ and every $x\in\T$. Hence, by Exercise~\ref{exer:expanding-maps-equiv-definition}, we conclude that every $g\in\mc{U}$ is an expanding map, and this concludes the proof.
\end{proof}


\subsubsection{Shadowing lemma for expanding endomorphisms}
\label{sec:shadowing-lemma-expandindg-circle}

Let $(X,d)$ be a compact metric space and $f\colon X\carr$ be a continuous map. Given real numbers $\alpha,\beta>0$ and a $\beta$-pseudo-orbit $(x_n)_{n\geq 0}$ for $f$, we say that the point $x\in X$ \textbf{$\alpha$-shadows} the pseudo-orbit $(x_n)_{n\geq 0}$ when it holds
\begin{displaymath}
  d\big(f^n(x),x_n\big)<\alpha, \quad\forall n\geq 0.
\end{displaymath}

The so-called \emph{shadowing property} plays a fundamental role in the theory of dynamical systems, specially in the proof of structural stability of hyperbolic systems, as well in the numerical study of them:

\begin{theorem}[Shadowing lemma for circle expanding maps]
  \label{thm:shadowing-lemma-circ-expanding}
  Let $f\in C_d^1(\T)$ be an expanding map. Then, there is a real number $\alpha_0>0$ such that for every $\alpha\in (0,\alpha_0)$, there is $\beta>0$ such that every $\beta$-pseudo-orbit $(x_n)_{n\geq 0}\subset\T$ for $f$, there is a unique point $x\in\T$ that $\alpha$-shadows the pseudo-orbit $(x_n)_{n\geq 0}$.
\end{theorem}

\begin{proof}
  Given $f$, let $C>0$ and $\lambda>1$ be such that $\abs{Df^n(x)}\geq C\lambda^n$, for every $x\in\T$ and $n\in\N$. Since $f$ is a $C^1$ map and it has no critical points and $\T$ is compact, there is $\alpha_0>0$ such that for every $x,y\in\T$, with $f(x)=y$, there is a local inverse $f_{x,y}^{-1}\colon B_\delta(y)\to\T$ of $f$, \ie~$f\circ f_{x,y}^{-1}(z)=z$, for every $z\in B_\delta(y)$. Then, we choose any $\alpha_0>0$ such that $\alpha_0<\min\{C(1-\lambda^{-1})\delta,\delta\}$. Given any $\alpha\in (0,\alpha_0)$, we shall determine $\beta>0$ afterwards as a function of $\alpha$. For the time being, let us just suppose $\beta<\alpha$. Then, given $\beta$-pseudo-orbit $(x_n)_{n\geq 0}$ for $f$, for each $m\in\N$ and each $n\in\{1,\ldots,m\}$ we consider the point
  \begin{displaymath}
    y_m^n:= f_{x_{m-n},x_{m-n+1}}^{-1}\circ\cdots\circ f_{x_{m-2},x_{m-1}}^{-1} \circ f_{x_{m-1},x_m}^{-1}(x_m), \quad\forall n\geq 1,
  \end{displaymath}
  and similarly we define $y_m^0:=x_m$, for all $m\in\N_0$.
  We claim that all these points are indeed all well defined. To see that, observe that for each $m\in\N$ it holds
  \begin{displaymath}
    \begin{split}
      d\big(y_m^1,y_m^0\big)&=d\big(f_{x_m,x_{m-1}}^{-1}(x_m), f_{x_m,x_{m-1}}^{-1}\circ f(x_{m-1})\big)\\
                              &\leq C^{-1}\lambda^{-1}d\big(x_m,f(x_{m-1})\big)<C^{-1}\lambda^{-1}\beta,
    \end{split}
  \end{displaymath}
  and by induction in $n\in\N$ we get
  \begin{equation}
    \label{eq:y-m-n-y-m-1-n-1-estimate}
    \begin{split}
      d\big(y_m^n,y_{m+1}^{n+1}\big) &= d\big(f_{x_{m-n},x_{m-n+1}}^{-1}(y_m^{n-1}), f_{x_{m-n},x_{m-n+1}}^{-1}(y_{m+1}^{n})\big)\\
                                     &\leq C^{-1}\lambda^{-n} d(y_m^0,y_m^1) < C^{-1}\lambda^{-n}\beta,
    \end{split}
  \end{equation}
  and hence
  \begin{equation}
    \label{eq:y-m-n-x-m-n-estimate}
    \begin{split}
      d(y_m^n,x_{m-n})  &\leq \sum_{j=1}^{n} d\big(y_{m-j+1}^{n-j+1},y_{m-j}^{n-j}\big)\\
                        &\leq  C^{-1}\alpha\sum_{j=1}^{n}\lambda^{-j} < \frac{C^{-1}\beta}{1-\lambda^{-1}}. \\
    \end{split}
  \end{equation}

  On the other hand, invoking estimate~\eqref{eq:y-m-n-y-m-1-n-1-estimate} we easily get that the sequence $(y_m^m)_{m\geq 0}$ is indeed a Cauchy sequence, and hence, it converges to some point $x\in\T$. Finally, by~\eqref{eq:y-m-n-x-m-n-estimate} we get that the point $x$ $\alpha$-shadows the pseudo-orbit $(x_n)_{n\geq 0}$ when $\beta<\alpha(1-\lambda^{-1})C$. The uniqueness of the point $x$ that $\alpha$-shadows the pseudo-orbit $(x_n)_{n\geq 0}$ easily follows by the fact that $f$ is an expanding map.
\end{proof}


\section{Dynamics in higher dimensions: local theory}
\label{sec:higher-dim-dynamics-local-theory}


In order to extend the theory of one-dimensional dynamics we developed in \S\ref{sec:real-line} and \S\ref{sec:dynamics-circle-maps} to higher dimensions, we first need to understand the dynamics of linear maps on finite dimensional vector spaces. More generally, it will be also useful to study the dynamics of bounded linear operators on Banach spaces.

\subsection{Linear operators on Banach spaces}\label{sec:line-oper-banach-spaces}

Let $(E,\norm{\cdot})$ be a Banach space over either $\C$ or $\R$. A linear operator $A\colon E\carr$ is said to be \textbf{bounded} when it is continuous and in such a case it holds
\begin{equation}\label{eq:norm-operator-def}
  \norm{A}:=\sup\{\norm{Ax} : x\in E,\ \norm{x}=1\}<\infty.
\end{equation}
We shall write $\mc{L}(E)$ for the vector space of bounded operators on $E$, which turns into a Banach algebra when it is endowed with norm given by~\eqref{eq:norm-operator-def}.

When $E$ is a Banach space over the complex field $\C$, we define the \textbf{spectrum} of the operator $A\in\mc{L}(E)$ as the set
\begin{displaymath}
  \sigma(A):=\left\{\lambda\in\C : A-\lambda I\ \text{is not bijective}\right\},
\end{displaymath}
where $I\in\mc{L}(E)$ denotes the identity operator.

When $E$ is a Banach space over the real field $\R$, one can define its \textbf{complexification} as the Banach space over $\C$ $(E_{\C},\norm{\cdot}_{\C})$ given by $E_{\C}:=E\times E$, $\norm{(x_{1},x_{2})}:=\max_{i}\norm{x_{i}}$, for every $(x_{1},x_{2})\in E_{\C}$ and the multiplication by a complex number is defined by
\begin{displaymath}
  (a+bi) (x_{1},x_{2}):=(ax_{1}-bx_{2},ax_{2}+bx_{1}), \quad\forall (x_{1},x_{2})\in E_{\C},\ \forall a+bi\in\C.
\end{displaymath}
We can define $E_{r}:=E\times\{0\}$ and $E_{i}:=\{0\}\times E$ as closed linear subspaces of $E_{\C}$ over $\R$, which are isomorphic to $E$ and such that $E_{\C}=E_{r}\oplus E_{i}$. Every operator $A\in\mc{L}(E)$ induces a new operator $A_{\C}\in\mc{L}(E_{\C})$ given by $A_{\C}(x_{1},x_{2}):=(Ax_{1},Ax_{2})$. On the other hand, an operator $B\in\mc{L}(E_{\C})$ is said to be a \textbf{real operator} when $B(E_{r})\subset E_{r}$ and $B(E_{i})\subset E_{i}$. One can show that every operator on $E_{\C}$ which is induced by another one on $E$ is indeed a real operator and reciprocally, every real operator on $E_{\C}$ is induced by one defined on $E$.

The \textbf{spectrum} of an operator $A\in\mc{L}(E)$ on the real Banach space $E$ is defined as the spectrum of the induced operator on the complexification of $E$, \ie~by $\sigma(A):=\sigma(A_{\C})$.

\begin{theorem}[Spectrum and spectral radius]\label{thm:spectrum-spectral-radius}
  If $E$ is a Banach space over either $\R$ or $\C$ and $A\in\mc{L}(E)\setminus\{0\}$, then it spectrum $\sigma(A)$ is a compact non-empty subset of $\C$. So, if we define the \textbf{spectral radius} of $A$ by
  \begin{displaymath}
    r(A):=\max\{\abs{\lambda} : \lambda\in\sigma(A)\},
  \end{displaymath}
  then it holds
  \begin{equation}\label{eq:formula-spectral-radius}
    r(A):=\lim_{n\to+\infty}\norm{A^{n}}^{1/n} = \inf_{n\geq 1}\norm{A^{n}}^{1/n}.
  \end{equation}

  On the other hand, if $A$ is an isomorphism, then it holds
  \begin{equation}\label{eq:spectrum-inverse}
    \sigma(A^{-1})=\left\{z^{-1}\in\C : z\in\sigma(A)\right\}.
  \end{equation}
\end{theorem}

\begin{proof}
  See for instance~\cite[\S1.2]{MurphyCAlgebrasOpTheory}.
\end{proof}

We shall write $\GL(E)$ for the group of invertible bounded linear maps on $E$, \ie~$\GL(E):=\{A\in\mc{L}(E) : 0\not\in\sigma(A)\}$.

\begin{definition}[Hyperbolic linear automorphisms]
  \label{def:hyperbolic-linear-automorphisms}
  Given a Banach space $E$, an invertible linear automorphism $A\in\GL(E)$ is said to be \textbf{hyperbolic} when it holds
  \begin{displaymath}
    \sigma(A)\cap\{\lambda\in\C : \abs{\lambda}=1\}=\emptyset.
  \end{displaymath}

  Since the spectrum $\sigma(A)$ is compact, when $A$ is hyperbolic there is constant $\lambda\in (0,1)$ such that
  \begin{equation}\label{eq:hyp-constant-def}
    \sigma(A)\cap\left\{z\in\C : \lambda\leq\abs{z}\leq\lambda^{-1}\right\}=\emptyset.
  \end{equation}
  Any number $\lambda\in(0,1)$ satisfying condition~\eqref{eq:hyp-constant-def} is called a \textbf{hyperbolicity constant} for $A$. The infimum of hyperbolicity constants shall be denoted by $A_h$
\end{definition}

The set of hyperbolic isomorphisms of $E$ will be denoted by $\mc{H}(E)$.

We will need the following result about spectral decompositions:
\begin{theorem}\label{thm:spectral-decomposition}
  Let $E$ be a Banach space and $A\in\mc{L}(E)$ be a bounded operator. Let $K_{1},K_{2}\subset\C$ be two non-empty compact sets such that $\sigma(A)=K_{1}\sqcup K_{2}$. Then there exists two closed proper subspaces $E_{1},E_{2}\subset E$ such that $E=E_{1}\oplus E_{2}$, $A(E_{i})\subset E_{i}$ and $\sigma\big(A\big|_{E_{i}}\big)=K_{i}$, for $i\in\{1,2\}$.
\end{theorem}

\begin{proof}
  See for instance~\cite[Chapter VII]{DunfordSchwartzLinOp}.
\end{proof}

As a consequence of Theorem~\ref{thm:spectrum-spectral-radius} we get the following result that characterizes the dynamics of hyperbolic linear maps:

\begin{theorem}[Hyperbolic linear maps]\label{thm:hyperbolic-linear-maps}
  Let $E$ be a Banach space and $A\in\mc{H}(E)$ be a hyperbolic linear automorphism. Let $E=E^{s}_{A}\oplus E^{u}_{A}$ be the linear splitting induced by the sets $K^{s}:=\sigma(A)\cap\{z\in\C : \abs{z}<1\}$ and $K^{u}:=\sigma(A)\cap\{z\in C : \abs{z}>1\}$ in Theorem~\ref{thm:spectral-decomposition}.

  Then, given any hyperbolicity constant $\lambda\in\big(0,1\big)$ for $A$ (\ie~condition~\eqref{eq:hyp-constant-def} holds), there exist a real constant $C>0$ such that
  \begin{displaymath}
    \norm{A^{n}v^{s}}\leq C\lambda^{n}\norm{v^{s}}\quad\text{and}\quad \norm{A^{-n}v^{u}}\leq C\lambda^{n}\norm{v^{u}},
  \end{displaymath}
  for all $n\in\N$ and every $v^{s}\in E^{s}_{A}$ and $v^{u}\in E^{u}_{A}$.

  Moreover, for every $\mu\in \big(A_h,1\big)$, where $A_h$ denotes the constant given by Definition~\ref{def:hyperbolic-linear-automorphisms}, there exists a norm $\abs{\cdot}$ on $E$, which is equivalent to $\norm{\cdot}$, and such that
  \begin{equation}\label{eq:adapted-norm-property}
    \abs{Av^{s}}\leq\mu\abs{v^{s}},\quad\text{and}\quad\abs{A^{-1}v^{u}}\leq\mu\abs{v^{u}},
  \end{equation}
  for every $v^{s}\in E^{s}_{A}$ and every $v^{u}\in E^{u}_{A}$.

  Any norm satisfying condition~\eqref{eq:adapted-norm-property} is an called an \textbf{adapted norm} for $A$.
\end{theorem}

\begin{corollary}[Dynamics of hyperbolic linear maps]\label{cor:dynamics-hyperbolic-linear-maps}
  If $E$ is a Banach space and $A\in\mc{H}(E)$, and $E=E^{s}_{A}\oplus E^{u}_{A}$ is the linear splitting given by Theorem~\ref{thm:hyperbolic-linear-maps}, then it holds
  \begin{align*}
    E^{s}_{A}&=\left\{x\in E : A^{n}x\to 0,\ \text{as } n\to+\infty\right\},\\
    E^{u}_{A}&=\left\{x\in E : A^{-n}x\to 0,\ \text{as } n\to+\infty\right\},
  \end{align*}
  and $A^{n}x\to\infty$, as $n\to\pm\infty$, for every $x\in E\setminus\big(E_{A}^{s}\cup E_{A}^{u})$.
\end{corollary}

\begin{proof}[Proof of Corollary~\ref{cor:dynamics-hyperbolic-linear-maps}]
  It is a straightforward consequence of Theorem~\ref{thm:hyperbolic-linear-maps}.
\end{proof}

\begin{proof}[Proof of Theorem~\ref{thm:hyperbolic-linear-maps}]
 Since $E^{s}_{A}$ and $E^{u}_{A}$ are closed subspaces, they are Banach spaces themselves. So, let us write $A_{i}:=A\big|_{E^{i}_{A}}\colon E^{i}_{A}\carr$, for $i\in\{s,u\}$. Let us fix a real number $\lambda\in \big(A_{h},1)$. By equation~\eqref{eq:formula-spectral-radius}, we know that there is $N\in\N$ such that $\norm{A^{n}_{s}}<\lambda^{n}$, and, by~\eqref{eq:spectrum-inverse}, $\norm{A_{u}^{-n}}<\lambda^{n}$, for every $n\geq N$.

 Let us write
 \begin{displaymath}
   C:=\lambda^{-N}\max_{1\leq n\leq N-1}\left\{\norm{A_{s}^{n}}, \norm{A_{u}^{-n}}\right\}.
 \end{displaymath}

 Then, given any $n\in\N$ there exist unique non-negative integers $r,q$, such that  $n=qN+r$ and $0\leq r<N$, and hence,
 \begin{displaymath}
   \norm{A_{s}^{n}}\leq \norm{A_{s}^{r}} \norm{A_{s}^{N}}^{q} \leq C\lambda^{N}\lambda^{Nq}< C\lambda^{Nq+r}=C\lambda^{n}.
 \end{displaymath}
 Similarly one gets
 \begin{displaymath}
   \norm{A_{u}^{-n}}\leq \norm{A_{s}^{-r}} \norm{A_{s}^{-N}}^{q} \leq C\lambda^{N}\lambda^{Nq}< C\lambda^{Nq+r}=C\lambda^{n}.
 \end{displaymath}

 In order to construct an adapted norm, let us fix a real number $\mu\in \big(A_{h},1)$ and let $\lambda$ be an arbitrary number in the open interval $\big(A_{H},\mu\big)$. By the above argument, there is a constant $C>0$ such that $\norm{A_{s}^{n}}\leq C\lambda^{n}$ and $\norm{A_{u}^{-n}}\leq C\lambda^{n}$, for every $n\in\N$. Then we define the new norm $\abs{\cdot}_{i}$ on the closed subspace $E_{A}^{i}$, for $i\in\{s,u\}$ by
 \begin{displaymath}
   \abs{x_{s}}_{s}:=\sum_{j=0}^{+\infty} \frac{1}{\mu^{j}}\norm{A^{j}x_{s}}, \quad\forall x_{s}\in E_{A}^{s},
 \end{displaymath}
 and
 \begin{displaymath}
   \abs{x_{u}}_{u}:=\sum_{j=0}^{+\infty} \frac{1}{\mu^{j}}\norm{A^{-j}x_{u}}, \quad\forall x_{u}\in E_{A}^{u}.
 \end{displaymath}

 Then, given an arbitrary vector $x\in E$, it can be uniquely written as $x=x_{s}+x_{u}$, with $x_{i}\in E_{A}^{i}$, an so we can define
 \begin{displaymath}
   \abs{x}:=\max\{\abs{x_{s}}_{s},\abs{x_{u}}_{u}\}.
 \end{displaymath}

 \begin{exercise}
   Show that $\abs{\cdot}_{s}$, $\abs{\cdot}_{u}$ and $\abs{\cdot}$ are indeed well define norms, and that $\abs{\cdot}$ and $\norm{\cdot}$ are equivalent, \ie~there exists a constant $K>1$ such that
   \begin{displaymath}
     K^{-1}\norm{x}\leq \abs{x}\leq K\norm{x}, \quad\forall x\in E.
   \end{displaymath}
 \end{exercise}

 Finally observe that
 \begin{displaymath}
   \abs{Ax_{s}} =\abs{Ax_{s}}_{s}=\sum_{j\geq 0} \frac{\norm{A^{j+1}x_{s}}}{\mu^{j}} = \mu (\abs{x_{s}}_{s}-\norm{x_{s}})\leq \mu\abs{x_{s}},
 \end{displaymath}
 for every $x_{s}\in E_{A}^{s}$ and
  \begin{displaymath}
   \abs{A^{-1}x_{u}} =\abs{A^{-1}x_{u}}_{u}=\sum_{j\geq 0} \frac{\norm{A^{-(j+1)}x_{s}}}{\mu^{j}} = \mu (\abs{x_{u}}_{u}-\norm{x_{u}})\leq \mu\abs{x_{u}},
 \end{displaymath}
 for every $x_{u}\in E^{u}_{A}$.
\end{proof}

\subsection{Non-linear hyperbolic local theory}
\label{sec:non-linear-local-theory}

There are three results that are fundamental in the local hyperbolic theory of non-linear dynamics, which are the following: the continuation of hyperbolic periodic points, Hartman-Grobman theorem and the stability of hyperbolic periodic orbits, and the stable/unstable manifold theorem.

We start showing the existence of continuation of hyperbolic points, which is a higher dimensional version of Exercise~\ref{ex:continuation-simple-periodic-points}:

\begin{theorem}
  \label{thm:cont-hyperbolic-fixed-point}
  Let $U\subset\R^d$ be an open subset, $f\colon U\to\R^d$ be a $C^1$ diffeomorphism onto its image and suppose that there is $p\in U$ such that $f(p)=p$ is a \textbf{simple fixed point,} \ie~$1$ does not belong to the spectrum of $Df_p$. Then, there exist $\eps>0$ and $r>0$ such that for every $C^1$ map $g\colon U\to\R^d$ which is $\eps$-$C^1$-close to $f$ (\ie $\abs{f(x)-g(x)}<\eps$ and $\norm{Df_x-Dg_x}<\eps$, for every $x\in\ U$), $\Fix(g)\cap B_\eps(p)$ is a singleton containing a unique point $p_g$, and $p_g$ depends continuously on $g$.
\end{theorem}

\begin{exercise}
  Prove Theorem~\ref{thm:cont-hyperbolic-fixed-point} using (the proof of) the implicit function theorem.
\end{exercise}

\bibliographystyle{amsalpha}
\bibliography{references}

\end{document}
